% Capítulo 10 (carpeta técnica) de Regiardo, Manuel. Estructura propia: ver PLAN.md del repositorio.
% Los archivos cu_01..04, requisitos, indice_cu, arquitectura y patrones se generan con scripts/cap10_gen.py
% a partir de documentation/legacy; no se editan a mano.
\capredactadotrue

% Figura a página completa con encabezado opcional (#1), que se coloca en la misma página que la figura.
% Si el diagrama es bastante más ancho que alto va en una página horizontal con los márgenes normales.
\newcommand{\figdiagh}[4]{%
  \sbox0{\includegraphics{#2}}%
  \ifdim\wd0>1.2\ht0
    \clearpage
    \begin{landscape}
      #1
      \begin{figure}[H]\centering
        \makebox[\textwidth]{\includegraphics[width=22.5cm,height=12.6cm,keepaspectratio]{#2}}
        \caption{#3}
      \end{figure}
    \end{landscape}
    \clearpage
  \else
    #1
    \begin{figure}[H]\centering
      \includegraphics[width=\textwidth,height=#4\textheight,keepaspectratio]{#2}
      \caption{#3}
    \end{figure}
  \fi}
\newcommand{\figdiag}[3]{\figdiagh{}{#1}{#2}{#3}}
\newcommand{\pendtxt}{\par\textit{[Pendiente de redacción]}\par}

\section{Aspectos Descriptivos de la Solución Tecnológica}

\subsection{Carátula}
\begin{tabularx}{\textwidth}{|l|X|}
\hline
\textbf{Proyecto} & \nombreproyecto \\ \hline
\textbf{Documento} & Carpeta técnica (Capítulo 10) \\ \hline
\textbf{Autor} & \autor{} (Legajo \legajo) \\ \hline
\textbf{Fecha} & \fechadoc \\ \hline
\end{tabularx}

\subsection{Índice}
El índice general de la carpeta técnica se encuentra al inicio del documento.

\subsection{Historial de revisión}
\begin{tabularx}{\textwidth}{|l|l|X|l|}
\hline
\textbf{Versión} & \textbf{Responsable} & \textbf{Descripción de la revisión} & \textbf{Fecha} \\ \hline
0.5 & Regiardo, Manuel & Versión preliminar de la carpeta técnica, elaborada a mano con algunos de los diagramas ya dibujados, presentada al profesor Sartorio, Alejandro & 04/08/2026 \\ \hline
1.0 & Regiardo, Manuel & Carpeta técnica con los casos de uso núcleo del DM & 17/11/2026 \\ \hline
\end{tabularx}

\subsection{Documento visión}
\subsubsection{Propuesta de valor del proyecto para la empresa}
\nombreproyecto{} centraliza en un único entorno digital la información de una campaña de juego de rol de mesa, sin reemplazar la experiencia presencial: mantiene el registro detallado de lo que ocurre alrededor de la partida para que el grupo dedique más tiempo a la narrativa y menos a la organización. Esta propuesta desarrolla, desde el punto de vista del sistema, la propuesta de valor para el cliente del Capítulo 1.

Para el \textbf{Dungeon Master (DM)}, el beneficio se concreta en cinco capacidades:
\begin{itemize}
  \item \textbf{Administrar campañas:} crear la campaña, convocar a los jugadores con un código de invitación y mantener el grupo sin planillas ni mensajes sueltos.
  \item \textbf{Gestionar sesiones:} acordar la fecha con el grupo mediante propuestas y votos, y llevar el estado de cada sesión (planificada, en curso, pausada o finalizada).
  \item \textbf{Personalizar las reglas:} adaptar el sistema de reglas de la campaña (por ejemplo, prohibir la magia o modificar un arma) sin afectar a otras campañas.
  \item \textbf{Generar diálogo de NPC:} obtener rápidamente el diálogo de un personaje improvisado, a partir de su raza, su clase y una idea aproximada, sin frenar la sesión.
  \item \textbf{Gestionar eventos:} registrar durante la partida el daño, la curación, la experiencia y el botín, y que los jugadores vean los cambios al instante y puedan consultar lo ocurrido.
\end{itemize}
Para el \textbf{jugador}, el valor está en unirse a la campaña, votar las fechas, conocer el estado de la sesión y ver reflejados al instante los cambios sobre su personaje.

\subsubsection{Especificaciones de requerimientos}
Los requisitos del sistema se organizan en cinco categorías: funcionales (RF), no funcionales (RNF), de dominio (RD), de sistema e interfaces (RS) y de entorno (RE). Se incluyen todos, aunque algunos no estén cubiertos por los casos de uso documentados en detalle en este documento. Al final se presentan las historias de usuario que los motivan.
% Generado por scripts/cap10_gen.py desde documentation/legacy. No editar a mano.
\paragraph{Requisitos Funcionales}

\subparagraph{Cuenta de usuario}
\begin{description}[leftmargin=2.2cm,labelwidth=2cm,itemsep=2pt]
\item[RF-001] Registrar un usuario nuevo (email, password, displayName) con cuenta propia, sin proveedor de identidad externo.
\item[RF-002] Iniciar sesión mediante email y password.
\item[RF-003] Roles globales de usuario (\texttt{user}, \texttt{admin}). Reservado para tareas de administración del sistema.
\item[RF-004] Cerrar sesión e invalidar credenciales de acceso.
\item[RF-005] Recuperar contraseña.
\end{description}
\subparagraph{Campañas}
\begin{description}[leftmargin=2.2cm,labelwidth=2cm,itemsep=2pt]
\item[RF-006] Un usuario autenticado crea una campaña y pasa a ser su DM.
\item[RF-007] El sistema genera un código de invitación único por campaña.
\item[RF-008] Un usuario se une a una campaña ingresando el código de invitación.
\item[RF-009] Un jugador abandona una campaña de la que forma parte.
\item[RF-010] El DM edita el nombre y la descripción de su campaña.
\item[RF-011] El DM elimina su campaña.
\item[RF-012] Un usuario lista las campañas en las que participa (como DM o Player).
\item[RF-013] El detalle de una campaña se devuelve filtrado según el rol del solicitante (el DM ve todo; el Player, una vista reducida).
\item[RF-014] El DM consulta y regenera el código de invitación de su campaña.
\item[RF-015] El DM remueve a un jugador de su campaña.
\item[RF-016] Al crear la campaña se asigna un sistema de reglas base (RuleSystem).
\end{description}
\subparagraph{Personajes}
\begin{description}[leftmargin=2.2cm,labelwidth=2cm,itemsep=2pt]
\item[RF-017] Un jugador crea un personaje en una campaña (raza, clase, nivel, estadísticas, HP, AC, habilidades, equipo, trasfondo).
\item[RF-018] Un usuario consulta el detalle de un personaje y la lista de sus personajes.
\item[RF-019] Un usuario actualiza o elimina sus propios personajes.
\item[RF-020] Los miembros de una campaña listan todos los personajes de la campaña.
\item[RF-021] Al crear o editar un personaje, el sistema ofrece el catálogo de razas, clases, estadísticas, traits, habilidades e ítems.
\end{description}
\subparagraph{Sesiones}
\begin{description}[leftmargin=2.2cm,labelwidth=2cm,itemsep=2pt]
\item[RF-022] El DM crea sesiones en una campaña (título, descripción, número de sesión). La fecha de inicio es opcional al crear.
\item[RF-023] Un jugador se une o abandona una sesión como participante.
\item[RF-024] El DM inicia, reanuda, pausa y finaliza una sesión (\texttt{planned $\rightarrow$ ongoing $\rightarrow$ paused/\allowbreak{}completed}).
\item[RF-025] Al iniciar o reanudar, el sistema compara el día de la acción con \texttt{start\_\allowbreak{}date}: mismo día, no modifica la fecha; día distinto, prevalece el momento real de inicio.
\item[RF-026] El DM edita las notas de la sesión (privadas y públicas).
\item[RF-027] El DM transfiere las notas (privadas, públicas o ambas) a otra sesión de la misma campaña; el sistema ofrece trasladar las de la sesión inmediatamente anterior al crear una sesión nueva.
\item[RF-028] El sistema oculta las notas privadas a los jugadores.
\item[RF-029] Un usuario consulta el listado de sesiones de una campaña y el detalle de una sesión.
\item[RF-030] El DM propone fechas para una sesión y quita propuestas mientras no haya una fecha confirmada.
\item[RF-031] Cada jugador vota sí, no o tal vez en cada fecha propuesta y puede cambiar su voto hasta la confirmación.
\item[RF-032] Los miembros ven los votos individuales y los totales por fecha; el sistema sugiere una fecha.
\item[RF-033] El DM confirma una fecha propuesta, que se guarda como fecha de inicio de la sesión.
\item[RF-034] El DM reprograma una sesión \texttt{planned} o \texttt{paused} fijando una nueva fecha de inicio.
\end{description}
\subparagraph{Eventos y bitácora}
\begin{description}[leftmargin=2.2cm,labelwidth=2cm,itemsep=2pt]
\item[RF-035] El DM registra eventos de sesión (daño, curación, XP, loot, personalizado) sobre uno, varios o todos los personajes de la campaña.
\item[RF-036] Al registrar un evento, el sistema lo difunde en tiempo real a los jugadores conectados a la sesión.
\item[RF-037] El DM redacta entradas de bitácora de tipo \texttt{combat} y \texttt{general}, opcionalmente asociadas a personajes.
\item[RF-038] Los miembros de la campaña consultan la bitácora completa de una sesión en orden cronológico.
\item[RF-039] Cada evento registrado genera automáticamente una entrada de bitácora de tipo \texttt{action}.
\item[RF-040] Los miembros de la campaña consultan el historial de eventos de una sesión.
\item[RF-041] Un cliente que se reconecta reconstruye el estado de la sesión consultando bitácora, eventos y fichas.
\end{description}
\subparagraph{Enemigos}
\begin{description}[leftmargin=2.2cm,labelwidth=2cm,itemsep=2pt]
\item[RF-042] El DM gestiona enemigos de la campaña (nombre, tipo, challenge rating, HP, AC, atributos, resistencias, acciones, loot).
\end{description}
\subparagraph{RuleSets}
\begin{description}[leftmargin=2.2cm,labelwidth=2cm,itemsep=2pt]
\item[RF-043] El DM modifica estadísticas de un ítem o hechizo dentro de su campaña (\texttt{modify\_\allowbreak{}stat}).
\item[RF-044] El DM banea un arma, hechizo o categoría (p. ej. magia) dentro de su campaña (\texttt{ban}).
\item[RF-045] El DM consulta el RuleSet efectivo de su campaña (reglas base más overrides).
\item[RF-046] El DM quita un override existente.
\item[RF-047] Los overrides rigen solo en la campaña donde se definieron.
\end{description}
\subparagraph{Diálogo de NPC}
\begin{description}[leftmargin=2.2cm,labelwidth=2cm,itemsep=2pt]
\item[RF-048] Durante una sesión en curso, el DM define los parámetros del NPC (raza, clase, rasgos de estado) y una idea del diálogo, y solicita su generación.
\item[RF-049] El DM regenera o descarta el diálogo generado.
\item[RF-050] El DM decide si guarda el diálogo como entrada de bitácora general; no se guarda automáticamente.
\item[RF-051] El DM consulta el historial de diálogos guardados.
\item[RF-052] Si el servicio de IA no está disponible, el sistema informa el error al DM y el resto de la aplicación sigue operando.
\end{description}
\subparagraph{Interfaz de usuario}
\begin{description}[leftmargin=2.2cm,labelwidth=2cm,itemsep=2pt]
\item[RF-053] Las rutas privadas exigen sesión activa y redirigen a \texttt{/\allowbreak{}login} si no la hay.
\item[RF-054] La aplicación ofrece páginas para login, registro, home, campañas, personajes, sesiones, bitácora y generación de diálogo de NPC.
\end{description}
\paragraph{Requisitos No Funcionales}
\begin{description}[leftmargin=2.2cm,labelwidth=2cm,itemsep=2pt]
\item[RNF-001] (Seguridad) Toda ruta protegida valida la credencial del usuario y resuelve el usuario de contexto antes de continuar. Aplica a todos los casos de uso salvo registro e inicio de sesión.
\item[RNF-002] (Autorización) Las operaciones sobre campañas y sesiones se validan según el rol del usuario en esa campaña (DM o Player), independientemente de su rol global. Aplica a todos los casos de uso de campañas, sesiones, eventos, RuleSets y diálogo de NPC.
\item[RNF-003] (Validación) Todo payload de entrada se valida contra un esquema antes de llegar al controlador.
\item[RNF-004] (Confidencialidad por rol) El filtrado de datos por rol se realiza en el servidor; un Player nunca recibe datos exclusivos del DM.
\item[RNF-005] (Seguridad de credenciales) Las contraseñas se almacenan con hash, nunca en texto plano.
\item[RNF-006] (Resiliencia) La caída o latencia de un servicio externo (API DnD5e, servicio de IA) no bloquea al resto del sistema; se espera manejo de errores y tiempos de espera.
\item[RNF-007] (Consistencia en tiempo real) La difusión por WebSocket es solo un canal de entrega; SQL Server es la única fuente de verdad.
\item[RNF-008] (Mantenibilidad) Todo el código está tipado con TypeScript.
\item[RNF-009] (Escalabilidad modular) El diseño modular permite agregar módulos sin afectar a los existentes.
\item[RNF-010] (Extensibilidad de reglas) Agregar un sistema de reglas nuevo no requiere modificar los existentes.
\item[RNF-011] (Usabilidad) Estilos consistentes y layout común para todas las vistas autenticadas.
\item[RNF-012] (Fechas) Las fechas se almacenan en UTC; la comparación de días usa la zona horaria del servidor y los clientes convierten a la zona del usuario.
\end{description}
\paragraph{Requisitos de Dominio}
\begin{description}[leftmargin=2.2cm,labelwidth=2cm,itemsep=2pt]
\item[RD-001] Los roles DM y Player son \textbf{por campaña}; un mismo usuario puede ser DM de una campaña y Player de otra.
\item[RD-002] Cada campaña tiene exactamente un DM, su creador.
\item[RD-003] El código de invitación es único por campaña.
\item[RD-004] Una sesión tiene estado \texttt{planned}, \texttt{ongoing}, \texttt{paused} o \texttt{completed}; transiciona \texttt{planned $\rightarrow$ ongoing $\rightarrow$ paused $\leftrightarrow$ ongoing $\rightarrow$ completed}.
\item[RD-005] La fecha de inicio de una sesión es opcional hasta que el DM confirma una fecha; la fecha confirmada es \texttt{start\_\allowbreak{}date} (no existe otro indicador de confirmación).
\item[RD-006] Un jugador tiene un único voto por propuesta (sí, no o tal vez) y puede cambiarlo hasta que se confirme la fecha; no hay fecha límite.
\item[RD-007] La fecha sugerida es la de más votos "sí"; ante empate, la de menos votos "no"; si persiste, no se sugiere ninguna. La sugerencia es orientativa y el DM decide.
\item[RD-008] Reprogramar no reabre la votación ni cambia el estado; una sesión \texttt{paused} reprogramada sigue \texttt{paused} hasta que el grupo se reúna y el DM la reanude.
\item[RD-009] Las notas de sesión son dos campos de texto plano; las privadas son solo del DM y las públicas son visibles para los jugadores.
\item[RD-010] Transferir notas es una copia por agregado entre sesiones de la misma campaña, sin vínculo ni sincronización posterior.
\item[RD-011] Los tipos de evento son \texttt{damage}, \texttt{healing}, \texttt{xp}, \texttt{loot} y \texttt{custom}.
\item[RD-012] Cada evento genera una entrada de bitácora \texttt{action}; los tipos de bitácora son \texttt{combat}, \texttt{general} y \texttt{action}.
\item[RD-013] Cada campaña tiene un único RuleSet, basado en un RuleSystem.
\item[RD-014] Un override es de tipo \texttt{ban} o \texttt{modify\_\allowbreak{}stat}, apunta a un ítem, hechizo, arma o categoría del catálogo base y es propio de la campaña.
\item[RD-015] Un diálogo de NPC solo puede generarse en una sesión en curso y no se persiste salvo que el DM lo guarde como entrada general.
\item[RD-016] Un personaje pertenece a una única campaña y a un único dueño; su nivel va de 1 a 20.
\item[RD-017] Los eventos de sesión solo pueden registrarse mientras la sesión está en curso.
\item[RD-018] Al aplicar eventos de daño y de curación, los puntos de vida del personaje se mantienen entre 0 y su máximo.
\item[RD-019] Un evento de loot se rechaza si alguno de sus ítems está baneado en el RuleSet de la campaña.
\item[RD-020] Todos los miembros de la campaña pueden ver todos los eventos de una sesión.
\end{description}
\paragraph{Requisitos de Sistema e Interfaces}
\begin{description}[leftmargin=2.2cm,labelwidth=2cm,itemsep=2pt]
\item[RS-001] Integración con la API pública DnD5e para obtener los catálogos del juego, con caché.
\item[RS-002] El diálogo de NPC lo genera un servicio externo independiente, consumido por HTTP/JSON desde el módulo \texttt{npc-dialogue}.
\item[RS-003] Persistencia relacional en SQL Server, con Prisma como ORM y migraciones, y tablas intermedias para relaciones muchos-a-muchos.
\item[RS-004] Canal de tiempo real por WebSocket (Socket.IO), un canal por sesión; el cliente envía su JWT en el handshake y el servidor verifica credencial y pertenencia a la campaña.
\item[RS-005] Autenticación propia con JWT.
\item[RS-006] Arquitectura de 3 capas cliente-servidor; backend estratificado y modular (\texttt{Routes $\rightarrow$ Controller $\rightarrow$ Service $\rightarrow$ Repository}); frontend MVC.
\item[RS-007] Los RuleSets se implementan con Strategy (un proveedor por sistema), Decorator (overrides por campaña) y Factory (selección del proveedor).
\item[RS-008] El servicio de diálogo es una decisión de arquitectura: un pipeline Pipes \& Filters con una arquitectura Blackboard en la construcción del prompt (ver \texttt{architecture.\allowbreak{}md}, sección 4). El proveedor de IA se oculta tras una Facade (patrón de diseño).
\item[RS-009] Frontend como SPA con React, Redux Toolkit (estado de autenticación global), react-router y TailwindCSS.
\item[RS-010] Backend en Node.js, TypeScript y Express.
\end{description}
\paragraph{Requisitos de Entorno}
\begin{description}[leftmargin=2.2cm,labelwidth=2cm,itemsep=2pt]
\item[RE-001] El backend y la base de datos se ejecutan en contenedores LXC sobre Proxmox VE; el artefacto de despliegue se construye con Dockerfile y Makefile.
\item[RE-002] El despliegue de demostración corre en un servidor propio (Proxmox VE) con NGINX como proxy inverso, que debe reenviar las cabeceras de actualización a WebSocket, y Cloudflare como proveedor de DNS.
\item[RE-003] SQL Server requiere al menos 2 GB de RAM.
\item[RE-004] Clientes: navegador web para el DM y aplicación móvil para los jugadores.
\item[RE-005] La configuración se entrega por variables de entorno; no se versionan secretos.
\end{description}
\paragraph{Historias de usuario}
\begin{description}[leftmargin=2.2cm,labelwidth=2cm,itemsep=3pt]
\item[HU-01] Como \textbf{DM} quiero crear una campaña y compartir un código para que mis jugadores se unan sin cargar sus datos a mano. \textit{(Requisitos: RF-006, RF-007, RF-008, RF-014)}
\item[HU-02] Como \textbf{DM} quiero elegir el sistema de reglas de mi campaña para que los personajes se creen con el juego que usamos. \textit{(Requisitos: RF-016, RF-021)}
\item[HU-03] Como \textbf{DM} quiero desactivar ítems o cambiar sus valores (p. ej. sin magia) para adaptar las reglas a mi mesa. \textit{(Requisitos: RF-043, RF-044, RF-045, RF-046, RF-047)}
\item[HU-04] Como \textbf{DM} quiero remover a un jugador o eliminar la campaña para mantener ordenado mi grupo. \textit{(Requisitos: RF-011, RF-015)}
\item[HU-05] Como \textbf{DM} quiero proponer fechas y ver quién puede cada una para coordinar el próximo encuentro del grupo. \textit{(Requisitos: RF-030, RF-032, RF-033)}
\item[HU-06] Como \textbf{Jugador} quiero votar las fechas propuestas y ver las respuestas de los demás para que nos pongamos de acuerdo. \textit{(Requisitos: RF-031, RF-032)}
\item[HU-07] Como \textbf{DM} quiero reprogramar una sesión cuando el grupo no puede o la sesión se extendió para continuar en el próximo encuentro. \textit{(Requisitos: RF-034, RF-024, RF-025)}
\item[HU-08] Como \textbf{DM} quiero iniciar, pausar y terminar una sesión para llevar el estado de la partida. \textit{(Requisitos: RF-024)}
\item[HU-09] Como \textbf{DM} quiero escribir notas privadas y públicas y llevarlas a la sesión siguiente para recordar lo que importa. \textit{(Requisitos: RF-026, RF-027, RF-028)}
\item[HU-10] Como \textbf{DM} quiero registrar eventos (daño, XP, loot) y que los jugadores los vean al instante para no frenar la partida. \textit{(Requisitos: RF-035, RF-036, RF-039)}
\item[HU-11] Como \textbf{Jugador} quiero ver el historial de lo ocurrido para retomar la historia si me reconecto. \textit{(Requisitos: RF-038, RF-040, RF-041)}
\item[HU-12] Como \textbf{DM} quiero generar el diálogo de un NPC improvisado indicando su raza, clase y una idea para no frenar la sesión. \textit{(Requisitos: RF-048, RF-049)}
\item[HU-13] Como \textbf{DM} quiero guardar los diálogos que me gustaron y consultarlos luego para reutilizarlos. \textit{(Requisitos: RF-050, RF-051)}
\item[HU-14] Como \textbf{DM} quiero que, si la IA falla, el resto de la sesión siga funcionando. \textit{(Requisitos: RF-052, RNF-006)}
\item[HU-15] Como \textbf{Jugador} quiero crear mi personaje dentro de las reglas de la campaña. \textit{(Requisitos: RF-017, RF-021, RF-047)}
\end{description}

\subsubsection{Alcance del desarrollo}
Esta carpeta técnica especifica en detalle los siguientes casos de uso núcleo, junto con los casos de uso que incluyen o extienden:
\begin{itemize}
  \item \textbf{CU-004-001} Gestionar Sesiones de Campaña, que comprende la planificación de la fecha con el grupo.
  \item \textbf{CU-007-001} Generar Diálogo de NPC.
  \item \textbf{CU-006-001} Personalizar RuleSet de Campaña.
  \item \textbf{CU-002-001} Administrar Campañas.
  \item \textbf{CU-005-001} Gestionar Eventos de Sesión, que comprende la bitácora de la sesión.
\end{itemize}
Los demás casos de uso del sistema se listan en el índice de casos de uso, pero no se especifican en detalle.

El alcance no incluye la gestión de enemigos, la recuperación de contraseña ni el soporte de sistemas de reglas distintos de D\&D 5e.

\subsubsection{Definiciones, acrónimos y abreviaciones}
\paragraph{Definiciones}
% Generado por scripts/cap10_gen.py desde documentation/legacy. No editar a mano.
{\small\setlength{\tabcolsep}{4pt}\renewcommand{\arraystretch}{1.15}
\begin{longtable}{|>{\raggedright\arraybackslash}p{4.2cm}|>{\raggedright\arraybackslash}p{11.1cm}|}
\hline
\textbf{Término} & \textbf{Definición} \\ \hline\endhead
\textbf{Campaña (Campaign)} & Partida de rol de largo plazo gestionada por un DM, con jugadores, personajes y sesiones asociadas. \\ \hline
\textbf{DM (Dungeon Master)} & Usuario que dirige la campaña; rol con permisos elevados dentro de esa campaña. \\ \hline
\textbf{Player} & Usuario miembro de una campaña que participa con uno o más personajes. \\ \hline
\textbf{Sesión (Session)} & Instancia de juego dentro de una campaña, con estado propio y notas. \\ \hline
\textbf{Personaje (Character)} & Ficha de personaje jugable perteneciente a un jugador dentro de una campaña. \\ \hline
\textbf{Enemigo (Enemy)} & Criatura/NPC antagonista gestionado por el DM dentro de una campaña. \\ \hline
\textbf{Evento (Event / SessionEvent)} & Registro de un suceso mecánico dentro de una sesión (daño, curación, XP, loot, custom) que afecta a personajes. \\ \hline
\textbf{RuleSystem} & Sistema de reglas base de un juego de rol de mesa (p. ej. "D\&D 5e"), independiente de cualquier campaña. \\ \hline
\textbf{CampaignRuleSet} & Instancia de un \texttt{RuleSystem} asociada a una campaña puntual, capaz de contener personalizaciones del DM. \\ \hline
\textbf{RuleOverride} & Personalización concreta (ban o modificación de estadística) aplicada sobre un \texttt{CampaignRuleSet}. \\ \hline
\textbf{Bitácora de Sesión (SessionLog)} & Entrada de historial narrativo y cronológico de una sesión (\texttt{combat}, \texttt{general} o \texttt{action}); las entradas \texttt{action} se generan automáticamente al registrar un \texttt{SessionEvent}. \\ \hline
\textbf{Diálogo de NPC} & Texto de diálogo generado para un NPC no planificado, mediante un servicio externo de IA. \\ \hline
\textbf{Propuesta de fecha (DateProposal)} & Fecha y horario candidatos que el DM propone para una sesión. \\ \hline
\textbf{Voto de fecha (DateVote)} & Respuesta de un jugador (sí / no / tal vez) a una propuesta de fecha. \\ \hline
\textbf{Notas de sesión} & Dos campos de texto plano por sesión: privadas (solo DM) y públicas (visibles para jugadores). \\ \hline
\end{longtable}}


\paragraph{Acrónimos y abreviaciones}
{\small\setlength{\tabcolsep}{4pt}
\begin{longtable}{|>{\raggedright\arraybackslash}p{3.2cm}|>{\raggedright\arraybackslash}p{12.0cm}|}
\hline
\textbf{Sigla} & \textbf{Significado} \\ \hline\endhead
DM & \textit{Dungeon Master}: director de juego de una campaña. \\ \hline
TTRPG & \textit{Tabletop Role-Playing Game}: juego de rol de mesa. \\ \hline
D\&D 5e & \textit{Dungeons \& Dragons}, quinta edición, sistema de reglas soportado. \\ \hline
NPC & \textit{Non-Player Character}: personaje no jugador, manejado por el DM. \\ \hline
CU & Caso de uso; se identifica con el formato CU-CCC-XXX. \\ \hline
RF, RNF, RD, RS, RE & Requisitos funcionales, no funcionales, de dominio, de sistema e interfaces y de entorno. \\ \hline
HU & Historia de usuario. \\ \hline
API & \textit{Application Programming Interface}. \\ \hline
REST & \textit{Representational State Transfer}: estilo de API sobre HTTP. \\ \hline
JWT & \textit{JSON Web Token}: credencial firmada para autenticar solicitudes. \\ \hline
ORM & \textit{Object-Relational Mapping}: capa de acceso a datos relacional (Prisma). \\ \hline
SPA & \textit{Single-Page Application}: aplicación web de una sola página. \\ \hline
MVC & Modelo-Vista-Controlador. \\ \hline
IA, LLM & Inteligencia artificial; \textit{Large Language Model}, modelo de lenguaje. \\ \hline
UML & \textit{Unified Modeling Language}. \\ \hline
ICONIX & Proceso de diseño guiado por casos de uso (especificación, robustez, secuencia). \\ \hline
UTC & Tiempo universal coordinado, en el que se almacenan las fechas. \\ \hline
\end{longtable}}
\subsubsection{Descripción del entorno del sistema}
\paragraph{Plataforma del lado del cliente}
\begin{itemize}
  \item \textbf{DM:} aplicación web de una sola página (React), que se usa desde un navegador moderno en una computadora de escritorio o portátil.
  \item \textbf{Jugador:} aplicación móvil, que se usa desde un teléfono inteligente.
  \item Ambos clientes consumen la misma API REST y reciben los eventos de sesión en tiempo real por WebSocket.
\end{itemize}

\paragraph{Plataforma del lado del servidor}
\begin{itemize}
  \item \textbf{Backend:} Node.js, TypeScript y Express, organizado en módulos por dominio con capas Rutas, Controlador, Servicio y Repositorio.
  \item \textbf{Base de datos:} SQL Server, con Prisma como ORM y sistema de migraciones.
  \item \textbf{Autenticación:} propia, basada en JWT.
  \item \textbf{Tiempo real:} canal WebSocket (Socket.IO) por sesión, con verificación de la credencial al conectarse.
  \item \textbf{Servicio de generación de diálogo de NPC:} servicio independiente, consumido por HTTP.
  \item \textbf{Servicio externo:} API pública de D\&D 5e para obtener los catálogos del juego.
\end{itemize}

\paragraph{Requisitos de entorno}
\begin{itemize}
  \item El backend y la base de datos se ejecutan en contenedores LXC sobre Proxmox VE; el artefacto de despliegue se construye con Dockerfile y Makefile. SQL Server requiere al menos 2~GB de memoria.
  \item El despliegue se realiza en un servidor propio con Proxmox VE, con NGINX como proxy inverso (que debe reenviar las cabeceras de actualización a WebSocket) y Cloudflare como proveedor de DNS.
  \item La configuración se entrega por variables de entorno; no se versionan secretos.
  \item \textbf{Seguridad:} toda ruta protegida valida la credencial del usuario y su rol en la campaña; las contraseñas se almacenan con \textit{hash}.
  \item \textbf{Disponibilidad y respaldo:} se establecen en el Capítulo 9 (Factibilidad técnica).
\end{itemize}
Estas condiciones se registran como requisitos de entorno RE-001 a RE-005.

\subsection{Descripción global del producto}
Los diagramas de esta sección se ajustarán para respetar la notación estándar propia de cada tipo de diagrama (por ejemplo, el diagrama de robustez).

\figdiagh{\subsubsection{Diagrama de dominio}
El diagrama de dominio muestra los conceptos del problema, sus atributos clave y sus relaciones, sin detalles técnicos de implementación.}{regiardo_manuel/domain-diagram}{Diagrama de dominio de DnDManager.}{0.8}

\subsubsection{Índice de casos de uso}
Cada caso de uso se identifica con el formato CU-CCC-XXX, donde CCC es el componente y XXX el caso de uso dentro del componente; el caso de uso núcleo de cada componente es el 001. Para cada uno se indican los requisitos que cubre.
% Generado por scripts/cap10_gen.py desde documentation/legacy. No editar a mano.
\paragraph{Componente 001 -- Cuenta}
{\footnotesize\setlength{\tabcolsep}{4pt}\renewcommand{\arraystretch}{1.15}
\begin{longtable}{|>{\raggedright\arraybackslash}p{1.9cm}|>{\raggedright\arraybackslash}p{4.4cm}|>{\raggedright\arraybackslash}p{1.7cm}|>{\raggedright\arraybackslash}p{2.5cm}|>{\raggedright\arraybackslash}p{4.0cm}|}
\hline
\textbf{Código} & \textbf{Nombre} & \textbf{Tipo} & \textbf{Actor primario} & \textbf{Requisitos} \\ \hline\endhead
CU-001-001 & Gestionar Cuenta de Usuario & Núcleo & Usuario & RF-001, RF-002 \\ \hline
CU-001-002 & Registrar Usuario & Incluido & Usuario no autenticado & RF-001, RNF-005 \\ \hline
CU-001-003 & Iniciar Sesión & Incluido & Usuario no autenticado & RF-002, RS-005, RNF-005 \\ \hline
CU-001-004 & Cerrar Sesión & Extiende & Usuario & RF-004 \\ \hline
CU-001-005 & Recuperar Contraseña & Extiende & Usuario & RF-005 \\ \hline
\end{longtable}}
\paragraph{Componente 002 -- Campañas}
{\footnotesize\setlength{\tabcolsep}{4pt}\renewcommand{\arraystretch}{1.15}
\begin{longtable}{|>{\raggedright\arraybackslash}p{1.9cm}|>{\raggedright\arraybackslash}p{4.4cm}|>{\raggedright\arraybackslash}p{1.7cm}|>{\raggedright\arraybackslash}p{2.5cm}|>{\raggedright\arraybackslash}p{4.0cm}|}
\hline
\textbf{Código} & \textbf{Nombre} & \textbf{Tipo} & \textbf{Actor primario} & \textbf{Requisitos} \\ \hline\endhead
CU-002-001 & Administrar Campañas & Núcleo & DM & RF-006, RF-012, RF-013 \\ \hline
CU-002-002 & Crear Campaña & Incluido & DM & RF-006, RF-007, RF-016, RD-002, RD-003 \\ \hline
CU-002-003 & Consultar Campañas y Detalle & Incluido & DM / Jugador & RF-012, RF-013, RNF-004 \\ \hline
CU-002-004 & Editar Campaña & Extiende & DM & RF-010 \\ \hline
CU-002-005 & Eliminar Campaña & Extiende & DM & RF-011 \\ \hline
CU-002-006 & Gestionar Código de Invitación & Extiende & DM & RF-014, RD-003 \\ \hline
CU-002-007 & Remover Jugador de la Campaña & Extiende & DM & RF-015 \\ \hline
CU-002-008 & Unirse a una Campaña & Extiende & Jugador & RF-008 \\ \hline
CU-002-009 & Abandonar una Campaña & Extiende & Jugador & RF-009 \\ \hline
\end{longtable}}
\paragraph{Componente 003 -- Personajes}
{\footnotesize\setlength{\tabcolsep}{4pt}\renewcommand{\arraystretch}{1.15}
\begin{longtable}{|>{\raggedright\arraybackslash}p{1.9cm}|>{\raggedright\arraybackslash}p{4.4cm}|>{\raggedright\arraybackslash}p{1.7cm}|>{\raggedright\arraybackslash}p{2.5cm}|>{\raggedright\arraybackslash}p{4.0cm}|}
\hline
\textbf{Código} & \textbf{Nombre} & \textbf{Tipo} & \textbf{Actor primario} & \textbf{Requisitos} \\ \hline\endhead
CU-003-001 & Gestionar Personajes de Usuario & Núcleo & Jugador & RF-017, RF-018 \\ \hline
CU-003-002 & Crear Personaje & Incluido & Jugador & RF-017, RF-021, RD-016, RS-001 \\ \hline
CU-003-003 & Editar Personaje & Extiende & Jugador & RF-019 \\ \hline
CU-003-004 & Eliminar Personaje & Extiende & Jugador & RF-019 \\ \hline
CU-003-005 & Consultar Personajes & Incluido & Jugador / DM & RF-018, RF-020 \\ \hline
CU-003-006 & Consultar Catálogo D\&D5e & Incluido & Sistema & RF-021, RS-001 \\ \hline
\end{longtable}}
\paragraph{Componente 004 -- Sesiones (con Planificación)}
{\footnotesize\setlength{\tabcolsep}{4pt}\renewcommand{\arraystretch}{1.15}
\begin{longtable}{|>{\raggedright\arraybackslash}p{1.9cm}|>{\raggedright\arraybackslash}p{4.4cm}|>{\raggedright\arraybackslash}p{1.7cm}|>{\raggedright\arraybackslash}p{2.5cm}|>{\raggedright\arraybackslash}p{4.0cm}|}
\hline
\textbf{Código} & \textbf{Nombre} & \textbf{Tipo} & \textbf{Actor primario} & \textbf{Requisitos} \\ \hline\endhead
CU-004-001 & Gestionar Sesiones de Campaña & Núcleo & DM & RF-022, RF-029 \\ \hline
CU-004-002 & Crear Sesión & Incluido & DM & RF-022, RF-027 \\ \hline
CU-004-003 & Consultar Sesiones y Detalle & Incluido & DM / Jugador & RF-029, RF-028, RNF-004 \\ \hline
CU-004-004 & Gestionar Ciclo de Vida de la Sesión & Extiende & DM & RF-024, RF-025, RD-004, RNF-012 \\ \hline
CU-004-005 & Proponer Fechas & Incluido & DM & RF-030 \\ \hline
CU-004-006 & Consultar Propuestas y Resultados & Incluido & DM / Jugador & RF-032, RD-007 \\ \hline
CU-004-007 & Confirmar Fecha & Incluido & DM & RF-033, RD-005 \\ \hline
CU-004-008 & Reprogramar Sesión & Extiende & DM & RF-034, RD-008 \\ \hline
CU-004-009 & Votar Fechas & Extiende & Jugador & RF-031, RD-006 \\ \hline
CU-004-010 & Unirse / Abandonar Sesión & Extiende & Jugador & RF-023 \\ \hline
\end{longtable}}
\paragraph{Componente 005 -- Eventos y Bitácora}
{\footnotesize\setlength{\tabcolsep}{4pt}\renewcommand{\arraystretch}{1.15}
\begin{longtable}{|>{\raggedright\arraybackslash}p{1.9cm}|>{\raggedright\arraybackslash}p{4.4cm}|>{\raggedright\arraybackslash}p{1.7cm}|>{\raggedright\arraybackslash}p{2.5cm}|>{\raggedright\arraybackslash}p{4.0cm}|}
\hline
\textbf{Código} & \textbf{Nombre} & \textbf{Tipo} & \textbf{Actor primario} & \textbf{Requisitos} \\ \hline\endhead
CU-005-001 & Gestionar Eventos de Sesión & Núcleo & DM & RF-035, RD-011, RD-017 \\ \hline
CU-005-002 & Difundir Evento en Tiempo Real & Incluido & Sistema & RF-036, RNF-007, RS-004 \\ \hline
CU-005-003 & Registrar Entrada de Acción en Bitácora & Incluido & Sistema & RF-039, RD-012 \\ \hline
CU-005-004 & Registrar Evento de Daño & Extiende & DM & RF-035, RD-011, RD-017, RD-018 \\ \hline
CU-005-005 & Registrar Evento de Curación & Extiende & DM & RF-035, RD-011, RD-017, RD-018 \\ \hline
CU-005-006 & Registrar Evento de XP & Extiende & DM & RF-035, RD-011, RD-017 \\ \hline
CU-005-007 & Registrar Evento de Loot & Extiende & DM & RF-035, RD-011, RD-017, RD-019 \\ \hline
CU-005-008 & Registrar Evento Personalizado & Extiende & DM & RF-035, RD-011, RD-017 \\ \hline
CU-005-009 & Consultar Historial de Eventos & Extiende & DM / Jugador & RF-040, RD-020 \\ \hline
CU-005-010 & Reconstruir Estado tras Reconexión & Extiende & Jugador / DM & RF-041, RNF-007 \\ \hline
CU-005-011 & Gestionar Bitácora de Sesión & Extiende & DM & RF-037, RF-038, RD-012 \\ \hline
CU-005-012 & Registrar Entrada de Combate & Incluido & DM & RF-037 \\ \hline
CU-005-013 & Registrar Entrada General & Incluido & DM & RF-037 \\ \hline
CU-005-014 & Consultar Bitácora & Incluido & DM / Jugador & RF-038 \\ \hline
\end{longtable}}
\paragraph{Componente 006 -- RuleSets}
{\footnotesize\setlength{\tabcolsep}{4pt}\renewcommand{\arraystretch}{1.15}
\begin{longtable}{|>{\raggedright\arraybackslash}p{1.9cm}|>{\raggedright\arraybackslash}p{4.4cm}|>{\raggedright\arraybackslash}p{1.7cm}|>{\raggedright\arraybackslash}p{2.5cm}|>{\raggedright\arraybackslash}p{4.0cm}|}
\hline
\textbf{Código} & \textbf{Nombre} & \textbf{Tipo} & \textbf{Actor primario} & \textbf{Requisitos} \\ \hline\endhead
CU-006-001 & Personalizar RuleSet de Campaña & Núcleo & DM & RF-043, RF-044, RF-047, RD-013, RD-014, RNF-010, RS-007 \\ \hline
CU-006-002 & Asignar Sistema de Reglas Base & Incluido & DM & RF-016, RD-013 \\ \hline
CU-006-003 & Consultar RuleSet Efectivo & Incluido & DM & RF-045 \\ \hline
CU-006-004 & Modificar Estadística de Ítem / Hechizo & Extiende & DM & RF-043, RD-014 \\ \hline
CU-006-005 & Banear Arma / Hechizo & Extiende & DM & RF-044, RD-014 \\ \hline
CU-006-006 & Quitar Override & Extiende & DM & RF-046 \\ \hline
\end{longtable}}
\paragraph{Componente 007 -- Diálogo de NPC}
{\footnotesize\setlength{\tabcolsep}{4pt}\renewcommand{\arraystretch}{1.15}
\begin{longtable}{|>{\raggedright\arraybackslash}p{1.9cm}|>{\raggedright\arraybackslash}p{4.4cm}|>{\raggedright\arraybackslash}p{1.7cm}|>{\raggedright\arraybackslash}p{2.5cm}|>{\raggedright\arraybackslash}p{4.0cm}|}
\hline
\textbf{Código} & \textbf{Nombre} & \textbf{Tipo} & \textbf{Actor primario} & \textbf{Requisitos} \\ \hline\endhead
CU-007-001 & Generar Diálogo de NPC & Núcleo & DM & RF-048, RD-015, RS-002, RNF-006 \\ \hline
CU-007-002 & Definir Parámetros del NPC & Incluido & DM & RF-048 \\ \hline
CU-007-003 & Construir Prompt (Blackboard) & Incluido & Servicio de IA & RS-008 \\ \hline
CU-007-004 & Solicitar Generación al Servicio de IA & Incluido & Servicio de IA & RS-002, RS-008 \\ \hline
CU-007-005 & Guardar Diálogo en Bitácora General & Extiende & DM & RF-050 \\ \hline
CU-007-006 & Regenerar Diálogo & Extiende & DM & RF-049 \\ \hline
CU-007-007 & Consultar Historial de Diálogos & Extiende & DM & RF-051 \\ \hline
CU-007-008 & Manejar Falla del Servicio de IA & Extiende & Sistema & RF-052, RNF-006 \\ \hline
\end{longtable}}
\paragraph{Componente 008 -- Notas}
{\footnotesize\setlength{\tabcolsep}{4pt}\renewcommand{\arraystretch}{1.15}
\begin{longtable}{|>{\raggedright\arraybackslash}p{1.9cm}|>{\raggedright\arraybackslash}p{4.4cm}|>{\raggedright\arraybackslash}p{1.7cm}|>{\raggedright\arraybackslash}p{2.5cm}|>{\raggedright\arraybackslash}p{4.0cm}|}
\hline
\textbf{Código} & \textbf{Nombre} & \textbf{Tipo} & \textbf{Actor primario} & \textbf{Requisitos} \\ \hline\endhead
CU-008-001 & Gestionar Notas de Sesión & Apoyo & DM & RF-026, RD-009 \\ \hline
CU-008-002 & Editar Notas de Sesión & Incluido & DM & RF-026 \\ \hline
CU-008-003 & Transferir Notas a Otra Sesión & Extiende & DM & RF-027, RD-010 \\ \hline
CU-008-004 & Consultar Notas de Sesión & Incluido & DM / Jugador & RF-028, RNF-004 \\ \hline
\end{longtable}}
\paragraph{Componente 009 -- Enemigos}
{\footnotesize\setlength{\tabcolsep}{4pt}\renewcommand{\arraystretch}{1.15}
\begin{longtable}{|>{\raggedright\arraybackslash}p{1.9cm}|>{\raggedright\arraybackslash}p{4.4cm}|>{\raggedright\arraybackslash}p{1.7cm}|>{\raggedright\arraybackslash}p{2.5cm}|>{\raggedright\arraybackslash}p{4.0cm}|}
\hline
\textbf{Código} & \textbf{Nombre} & \textbf{Tipo} & \textbf{Actor primario} & \textbf{Requisitos} \\ \hline\endhead
CU-009-001 & Gestionar Enemigos de Campaña & Núcleo & DM & RF-042 \\ \hline
CU-009-002 a CU-009-005 & Crear, Editar, Eliminar y Consultar Enemigos & Incluido / Extiende & DM & RF-042 \\ \hline
\end{longtable}}
\paragraph{Trazabilidad de requisitos}
Los requisitos \textbf{globales} se aplican a todos los casos de uso y no se repiten fila por fila: RNF-001 (autenticación; no aplica a registro e inicio de sesión), RNF-002 (autorización por rol de campaña), RNF-003 (validación de datos), RNF-008, RNF-009, RNF-011, RF-053, RF-054, RD-001, RS-003, RS-006, RS-009, RS-010 y RE-001 a RE-005.

{\small\setlength{\tabcolsep}{4pt}\renewcommand{\arraystretch}{1.15}
\begin{longtable}{|>{\raggedright\arraybackslash}p{2.4cm}|>{\raggedright\arraybackslash}p{12.9cm}|}
\hline
\textbf{Requisito} & \textbf{Casos de uso que lo cubren} \\ \hline\endhead
RF-001 & CU-001-001, CU-001-002 \\ \hline
RF-002 & CU-001-001, CU-001-003 \\ \hline
RF-003 & — \\ \hline
RF-004 & CU-001-004 \\ \hline
RF-005 & CU-001-005 \\ \hline
RF-006 & CU-002-001, CU-002-002 \\ \hline
RF-007 & CU-002-002 \\ \hline
RF-008 & CU-002-008 \\ \hline
RF-009 & CU-002-009 \\ \hline
RF-010 & CU-002-004 \\ \hline
RF-011 & CU-002-005 \\ \hline
RF-012 & CU-002-001, CU-002-003 \\ \hline
RF-013 & CU-002-001, CU-002-003 \\ \hline
RF-014 & CU-002-006 \\ \hline
RF-015 & CU-002-007 \\ \hline
RF-016 & CU-002-002, CU-006-002 \\ \hline
RF-017 & CU-003-001, CU-003-002 \\ \hline
RF-018 & CU-003-001, CU-003-005 \\ \hline
RF-019 & CU-003-003, CU-003-004 \\ \hline
RF-020 & CU-003-005 \\ \hline
RF-021 & CU-003-002, CU-003-006 \\ \hline
RF-022 & CU-004-001, CU-004-002 \\ \hline
RF-023 & CU-004-010 \\ \hline
RF-024 & CU-004-004 \\ \hline
RF-025 & CU-004-004 \\ \hline
RF-026 & CU-008-001, CU-008-002 \\ \hline
RF-027 & CU-004-002, CU-008-003 \\ \hline
RF-028 & CU-004-003, CU-008-004 \\ \hline
RF-029 & CU-004-001, CU-004-003 \\ \hline
RF-030 & CU-004-005 \\ \hline
RF-031 & CU-004-009 \\ \hline
RF-032 & CU-004-006 \\ \hline
RF-033 & CU-004-007 \\ \hline
RF-034 & CU-004-008 \\ \hline
RF-035 & CU-005-001, CU-005-004, CU-005-005, CU-005-006, CU-005-007, CU-005-008 \\ \hline
RF-036 & CU-005-002 \\ \hline
RF-037 & CU-005-011, CU-005-012, CU-005-013 \\ \hline
RF-038 & CU-005-011, CU-005-014 \\ \hline
RF-039 & CU-005-003 \\ \hline
RF-040 & CU-005-009 \\ \hline
RF-041 & CU-005-010 \\ \hline
RF-042 & CU-009-001, CU-009-002 a 009-005 \\ \hline
RF-043 & CU-006-001, CU-006-004 \\ \hline
RF-044 & CU-006-001, CU-006-005 \\ \hline
RF-045 & CU-006-003 \\ \hline
RF-046 & CU-006-006 \\ \hline
RF-047 & CU-006-001 \\ \hline
RF-048 & CU-007-001, CU-007-002 \\ \hline
RF-049 & CU-007-006 \\ \hline
RF-050 & CU-007-005 \\ \hline
RF-051 & CU-007-007 \\ \hline
RF-052 & CU-007-008 \\ \hline
RNF-004 & CU-002-003, CU-004-003, CU-008-004 \\ \hline
RNF-005 & CU-001-002, CU-001-003 \\ \hline
RNF-006 & CU-007-001, CU-007-008 \\ \hline
RNF-007 & CU-005-002, CU-005-010 \\ \hline
RNF-010 & CU-006-001 \\ \hline
RNF-012 & CU-004-004 \\ \hline
RD-002 & CU-002-002 \\ \hline
RD-003 & CU-002-002, CU-002-006 \\ \hline
RD-004 & CU-004-004 \\ \hline
RD-005 & CU-004-007 \\ \hline
RD-006 & CU-004-009 \\ \hline
RD-007 & CU-004-006 \\ \hline
RD-008 & CU-004-008 \\ \hline
RD-009 & CU-008-001 \\ \hline
RD-010 & CU-008-003 \\ \hline
RD-011 & CU-005-001, CU-005-004, CU-005-005, CU-005-006, CU-005-007, CU-005-008 \\ \hline
RD-012 & CU-005-003, CU-005-011 \\ \hline
RD-013 & CU-006-001, CU-006-002 \\ \hline
RD-014 & CU-006-001, CU-006-004, CU-006-005 \\ \hline
RD-015 & CU-007-001 \\ \hline
RD-016 & CU-003-002 \\ \hline
RD-017 & CU-005-001, CU-005-004, CU-005-005, CU-005-006, CU-005-007, CU-005-008 \\ \hline
RD-018 & CU-005-004, CU-005-005 \\ \hline
RD-019 & CU-005-007 \\ \hline
RD-020 & CU-005-009 \\ \hline
RS-001 & CU-003-002, CU-003-006 \\ \hline
RS-002 & CU-007-001, CU-007-004 \\ \hline
RS-004 & CU-005-002 \\ \hline
RS-005 & CU-001-003 \\ \hline
RS-007 & CU-006-001 \\ \hline
RS-008 & CU-007-003, CU-007-004 \\ \hline
\end{longtable}}


\subsubsection{Casos de uso}
Se documentan en detalle los casos de uso núcleo del Dungeon Master. Cada uno incluye su especificación, el prototipo de interfaz y los diagramas de casos de uso, de robustez y de secuencia.
% Generado por scripts/cap10_gen.py desde documentation/legacy. No editar a mano.
\paragraph{CU-004-001 -- Gestionar Sesiones de Campaña}
{\small\setlength{\tabcolsep}{5pt}\renewcommand{\arraystretch}{1.2}
\begin{longtable}{|>{\raggedright\arraybackslash\bfseries}p{3.2cm}|>{\raggedright\arraybackslash}p{11.9cm}|}
\hline
Objetivo & Permitir al DM planificar la fecha de una sesión de su campaña junto con el grupo y conducirla (iniciar, pausar, reanudar, finalizar), manteniendo a los jugadores informados de su estado y fecha. \\ \hline
Identificación & CU-004-001 \\ \hline
Nombre & Gestionar Sesiones de Campaña \\ \hline
Descripción & El DM crea una sesión dentro de su campaña, propone fechas candidatas y los jugadores votan sí, no o tal vez en cada una. El sistema muestra los votos de todos y sugiere una fecha, pero el DM decide y confirma. Llegado el momento, el DM inicia la sesión, puede pausarla y reanudarla en un encuentro posterior, y finalmente la finaliza. Si el grupo no puede reunirse, el DM reprograma la sesión fijando una nueva fecha. Al crear una sesión nueva, el sistema ofrece trasladar las notas de la sesión anterior. \\ \hline
Precondición & \begin{itemize}[itemsep=2pt,topsep=2pt]
\item El usuario está autenticado y es DM de la campaña.
\item La campaña tiene al menos un jugador para que la votación tenga sentido (no es obligatorio para crear la sesión).
\end{itemize} \\ \hline
Postcondición & \begin{itemize}[itemsep=2pt,topsep=2pt]
\item La sesión existe con su fecha de inicio confirmada (o sin fecha, si el DM aún no confirmó una) y en el estado que corresponde a la última acción del DM (\texttt{planned}, \texttt{ongoing}, \texttt{paused} o \texttt{completed}).
\end{itemize} \\ \hline
Actores primarios & DM \\ \hline
Actores secundarios & Jugador (vota fechas, consulta la sesión y sus propuestas). \\ \hline
Eventos disparadores & El DM decide crear una nueva sesión o gestionar una existente de su campaña. \\ \hline
Extensiones y condiciones & \begin{itemize}[itemsep=2pt,topsep=2pt]
\item \textbf{include}: CU-004-002 Crear Sesión. Siempre al planificar una sesión nueva
\item \textbf{include}: CU-004-003 Consultar Sesiones y Detalle. Siempre que se revisa el estado
\item \textbf{include}: CU-004-005 Proponer Fechas. Al planificar la fecha
\item \textbf{include}: CU-004-006 Consultar Propuestas y Resultados. Antes de confirmar
\item \textbf{include}: CU-004-007 Confirmar Fecha. Para fijar la fecha
\item \textbf{extend}: CU-004-004 Gestionar Ciclo de Vida. Al comenzar, pausar, reanudar o finalizar
\item \textbf{extend}: CU-004-008 Reprogramar Sesión. Sesión \texttt{planned} o \texttt{paused}
\item \textbf{extend}: CU-004-009 Votar Fechas. Lo realiza el Jugador
\item \textbf{extend}: CU-004-010 Unirse / Abandonar Sesión. Lo realiza el Jugador
\item \textbf{extend}: CU-008-001 Gestionar Notas de Sesión. Notas privadas y públicas; traslado entre sesiones
\end{itemize}
Requisitos relacionados: RF-022 a RF-034, RD-004 a RD-010, RNF-001, RNF-002, RNF-004, RNF-012 (ver HU-05 a HU-09). \\ \hline
Escenario principal & \begin{enumerate}[itemsep=2pt,topsep=2pt]
\item \textbf{DM:} solicita crear una sesión e ingresa título, descripción y número de sesión.
\item \textbf{Sistema:} crea la sesión en estado "planificada" sin fecha de inicio y ofrece trasladar las notas de la sesión anterior.
\item \textbf{DM:} acepta o rechaza el traslado de notas.
\item \textbf{DM:} propone una o más fechas candidatas.
\item \textbf{Sistema:} registra las propuestas y las muestra a los jugadores.
\item \textbf{Jugador:} responde sí, no o tal vez para cada fecha.
\item \textbf{Sistema:} registra los votos y muestra a todos los miembros los votos individuales y los totales por fecha.
\item \textbf{DM:} consulta los resultados; el sistema indica la fecha sugerida (más "sí"; ante empate, menos "no").
\item \textbf{DM:} confirma una fecha.
\item \textbf{Sistema:} fija la fecha de inicio de la sesión y cierra la votación.
\item \textbf{DM:} inicia la sesión el día acordado.
\item \textbf{Sistema:} marca la sesión como "en curso" y conserva la fecha de inicio.
\item \textbf{DM:} finaliza la sesión.
\item \textbf{Sistema:} marca la sesión como "completada".
\end{enumerate} \\ \hline
Escenarios alternativos & \begin{itemize}[itemsep=2pt,topsep=2pt]
\item \textbf{Paso 1 — sin fecha prevista:} el DM no propone fechas; la sesión queda planificada sin fecha hasta que confirme o reprograme una.
\item \textbf{Paso 2 — es la primera sesión:} no hay sesión anterior; el sistema no ofrece traslado de notas.
\item \textbf{Paso 3 — traslado de notas:} el DM elige qué trasladar (privadas, públicas o ambas); el sistema agrega ese texto al final de las notas de la nueva sesión, sin vínculo posterior.
\item \textbf{Paso 4 — quitar propuesta:} el DM elimina una propuesta; el sistema elimina también sus votos.
\item \textbf{Paso 6 — cambio de voto:} el jugador cambia su respuesta mientras no haya fecha confirmada.
\item \textbf{Paso 6 — votar con fecha confirmada:} el sistema rechaza el voto e informa que la votación está cerrada.
\item \textbf{Paso 8 — empate persistente:} el sistema no sugiere una fecha y el DM decide libremente.
\item \textbf{Paso 11 — día distinto al acordado:} el sistema registra como fecha de inicio el momento real en que se inició la sesión.
\item \textbf{Después del paso 11 — pausar:} el DM pausa la sesión; el sistema la marca como "pausada" y el DM la reanuda en un encuentro posterior.
\item \textbf{Reprogramar:} con la sesión planificada o pausada, el DM fija una nueva fecha; el sistema la actualiza sin reabrir la votación ni cambiar el estado (una sesión pausada sigue pausada).
\item \textbf{Usuario sin permiso:} si el usuario no es DM de la campaña, el sistema rechaza la operación.
\end{itemize} \\ \hline
\end{longtable}}
\subparagraph{Prototipo de la interfaz de usuario}
Pantallas: Crear Sesión; Planificar Fechas (propuestas, votos por jugador, fecha sugerida, botón Confirmar); Detalle de Sesión (estado, fecha, notas, botones Iniciar / Pausar / Finalizar / Reprogramar).

\figdiagh{\subparagraph{Diagrama de casos de uso}}{regiardo_manuel/uc-cu-004-001}{Diagrama de casos de uso de CU-004-001.}{0.9}
\figdiagh{\subparagraph{Diagrama de robustez}}{regiardo_manuel/robustness-cu-004-001}{Diagrama de robustez de CU-004-001.}{0.9}
\figdiagh{\subparagraph{Diagrama de secuencia}}{regiardo_manuel/sequence-cu-004-001}{Diagrama de secuencia de CU-004-001.}{0.9}

% Generado por scripts/cap10_gen.py desde documentation/legacy. No editar a mano.
\paragraph{CU-007-001 -- Generar Diálogo de NPC}
{\small\setlength{\tabcolsep}{5pt}\renewcommand{\arraystretch}{1.2}
\begin{longtable}{|>{\raggedright\arraybackslash\bfseries}p{3.2cm}|>{\raggedright\arraybackslash}p{11.9cm}|}
\hline
Objetivo & Permitir al DM obtener, durante una sesión en curso, un diálogo para un NPC no planificado a partir de pocos datos, sin frenar la partida y sin que una falla del servicio de IA afecte al resto de la aplicación. \\ \hline
Identificación & CU-007-001 \\ \hline
Nombre & Generar Diálogo de NPC \\ \hline
Descripción & Durante una sesión en curso, el DM indica la raza y la clase del NPC que habla, algunos rasgos de su estado y una idea aproximada de lo que debe decir. El sistema envía esos datos, junto con el contexto de la campaña, a un servicio de generación de diálogo, que arma el contexto completo (mundo, emociones, rasgos de habla) y devuelve el texto. El DM lo lee y decide: regenerarlo, descartarlo o guardarlo como entrada general de la bitácora de la sesión. El diálogo no se guarda automáticamente. Si el servicio no responde, el sistema lo informa y el DM sigue usando la aplicación. \\ \hline
Precondición & \begin{itemize}[itemsep=2pt,topsep=2pt]
\item El usuario está autenticado y es DM de la campaña.
\item La sesión está en curso (\texttt{ongoing}).
\end{itemize} \\ \hline
Postcondición & \begin{itemize}[itemsep=2pt,topsep=2pt]
\item El DM tiene el diálogo generado en pantalla, o un aviso de que el servicio no está disponible.
\item Si el DM lo guardó, existe una entrada general nueva en la bitácora de la sesión; en caso contrario, no se persiste nada.
\end{itemize} \\ \hline
Actores primarios & DM \\ \hline
Actores secundarios & Servicio de Generación de Diálogo (IA), servicio externo propio. \\ \hline
Eventos disparadores & El DM necesita diálogo para un NPC no planificado durante la sesión. \\ \hline
Extensiones y condiciones & \begin{itemize}[itemsep=2pt,topsep=2pt]
\item \textbf{extend (de CU-004-001)}: Gestionar Sesiones de Campaña. Requiere sesión en curso
\item \textbf{include}: CU-007-002 Definir Parámetros del NPC. Siempre, antes de generar
\item \textbf{include}: CU-007-003 Construir Prompt (Blackboard). Lo realiza el servicio externo
\item \textbf{include}: CU-007-004 Solicitar Generación al Servicio de IA. Siempre
\item \textbf{extend}: CU-007-005 Guardar Diálogo en Bitácora General. Solo si el DM lo decide
\item \textbf{extend}: CU-007-006 Regenerar Diálogo. Si el resultado no conforma al DM
\item \textbf{extend}: CU-007-007 Consultar Historial de Diálogos. Lista los diálogos guardados
\item \textbf{extend (de CU-007-004)}: CU-007-008 Manejar Falla del Servicio de IA. Si el servicio no responde o falla
\end{itemize}
Requisitos relacionados: RF-048 a RF-052, RD-015, RS-002, RS-008, RNF-006, RNF-001, RNF-002 (ver HU-12 a HU-14). \\ \hline
Escenario principal & \begin{enumerate}[itemsep=2pt,topsep=2pt]
\item \textbf{DM:} solicita generar un diálogo de NPC.
\item \textbf{Sistema:} muestra el formulario de parámetros.
\item \textbf{DM:} ingresa la raza y la clase del NPC, los rasgos de estado y una idea aproximada del diálogo.
\item \textbf{Sistema:} valida los datos y que la sesión esté en curso, y muestra que la generación está en proceso.
\item \textbf{Sistema:} muestra el diálogo generado.
\item \textbf{DM:} lee el diálogo y elige guardarlo en la bitácora.
\item \textbf{Sistema:} guarda el diálogo como entrada general de la bitácora de la sesión y lo confirma.
\end{enumerate} \\ \hline
Escenarios alternativos & \begin{itemize}[itemsep=2pt,topsep=2pt]
\item \textbf{Paso 4 — datos incompletos:} el sistema indica los campos faltantes y no solicita la generación.
\item \textbf{Paso 4 — sesión no en curso:} el sistema informa que la sesión debe estar en curso y no genera el diálogo.
\item \textbf{Paso 5 — el servicio de IA no responde o falla:} el sistema informa que la generación no está disponible; el DM puede reintentar y el resto de la aplicación sigue funcionando.
\item \textbf{Paso 6 — regenerar:} el DM pide otro diálogo; el sistema genera uno nuevo con los mismos parámetros y descarta el anterior.
\item \textbf{Paso 6 — descartar:} el DM no guarda el diálogo; el sistema no persiste nada.
\item \textbf{Paso 6 — cambiar parámetros:} el DM vuelve al formulario, modifica los datos y solicita una nueva generación.
\item \textbf{Después del paso 7 — historial:} el DM abre el historial y el sistema lista los diálogos guardados en la bitácora general de la sesión.
\item \textbf{Usuario sin permiso:} si el usuario no es DM de la campaña, el sistema rechaza la operación.
\end{itemize} \\ \hline
\end{longtable}}
\subparagraph{Prototipo de la interfaz de usuario}
Pantallas: Formulario de parámetros (raza, clase, rasgos de estado, idea); Resultado (texto generado y botones Guardar, Regenerar, Descartar); Historial de diálogos; mensaje de servicio no disponible.

\figdiagh{\subparagraph{Diagrama de casos de uso}}{regiardo_manuel/uc-cu-007-001}{Diagrama de casos de uso de CU-007-001.}{0.9}
\figdiagh{\subparagraph{Diagrama de robustez}}{regiardo_manuel/robustness-cu-007-001}{Diagrama de robustez de CU-007-001.}{0.9}
\figdiagh{\subparagraph{Diagrama de secuencia}}{regiardo_manuel/sequence-cu-007-001}{Diagrama de secuencia de CU-007-001.}{0.9}

\iftfi
% Generado por scripts/cap10_gen.py desde documentation/legacy. No editar a mano.
\paragraph{CU-006-001 -- Personalizar RuleSet de Campaña}
{\small\setlength{\tabcolsep}{5pt}\renewcommand{\arraystretch}{1.2}
\begin{longtable}{|>{\raggedright\arraybackslash\bfseries}p{3.2cm}|>{\raggedright\arraybackslash}p{11.9cm}|}
\hline
Objetivo & Permitir al DM adaptar las reglas de su campaña (por ejemplo, prohibir la magia o cambiar los valores de un arma) sin alterar el catálogo base ni las demás campañas. \\ \hline
Identificación & CU-006-001 \\ \hline
Nombre & Personalizar RuleSet de Campaña \\ \hline
Descripción & Cada campaña usa un sistema de reglas base (por ejemplo, D\&D 5e), elegido al crearla. Sobre esa base, el DM puede consultar las reglas que rigen en su campaña, modificar las estadísticas de un ítem o hechizo, banear un arma, un hechizo o una categoría entera (por ejemplo, la magia) y quitar una personalización cuando ya no la quiera. Las personalizaciones rigen solo en esa campaña y se aplican al crear personajes y repartir botín. El sistema está preparado para incorporar otros sistemas de reglas sin modificar los existentes. \\ \hline
Precondición & \begin{itemize}[itemsep=2pt,topsep=2pt]
\item El usuario está autenticado y es DM de la campaña.
\item La campaña tiene un sistema de reglas base asignado.
\end{itemize} \\ \hline
Postcondición & \begin{itemize}[itemsep=2pt,topsep=2pt]
\item El RuleSet de la campaña refleja los overrides vigentes: las reglas base con las personalizaciones aplicadas.
\item El catálogo base y las demás campañas no se modificaron.
\end{itemize} \\ \hline
Actores primarios & DM \\ \hline
Actores secundarios & API DnD5e (sistema de reglas externo que provee el catálogo base). \\ \hline
Eventos disparadores & El DM quiere ver o cambiar las reglas que rigen en su campaña. \\ \hline
Extensiones y condiciones & \begin{itemize}[itemsep=2pt,topsep=2pt]
\item \textbf{include}: CU-006-002 Asignar Sistema de Reglas Base. Al crear la campaña (incluido también por CU-002-002)
\item \textbf{include}: CU-006-003 Consultar RuleSet Efectivo. Siempre que se revisan las reglas
\item \textbf{extend}: CU-006-004 Modificar Estadística de Ítem / Hechizo. El DM define nuevos valores para un ítem
\item \textbf{extend}: CU-006-005 Banear Arma / Hechizo / Categoría. El DM marca algo como no disponible
\item \textbf{extend}: CU-006-006 Quitar Override. El DM revierte una personalización
\end{itemize}
Reglas: existe un único override por objetivo en cada campaña; el ítem a personalizar debe existir en el catálogo base.

Requisitos relacionados: RF-016, RF-043 a RF-047, RD-013, RD-014, RNF-010, RS-001, RS-007, RNF-001, RNF-002 (ver HU-02, HU-03). \\ \hline
Escenario principal & \begin{enumerate}[itemsep=2pt,topsep=2pt]
\item \textbf{DM:} abre el RuleSet de su campaña.
\item \textbf{Sistema:} muestra las reglas vigentes: el catálogo base con los overrides ya aplicados.
\item \textbf{DM:} elige un ítem del catálogo y define nuevos valores para sus estadísticas.
\item \textbf{Sistema:} confirma que el ítem existe y registra la modificación.
\item \textbf{DM:} marca un arma, un hechizo o una categoría como no disponible.
\item \textbf{Sistema:} registra el baneo.
\item \textbf{DM:} consulta nuevamente el RuleSet.
\item \textbf{Sistema:} muestra las reglas con las modificaciones y los baneos vigentes.
\end{enumerate} \\ \hline
Escenarios alternativos & \begin{itemize}[itemsep=2pt,topsep=2pt]
\item \textbf{Paso 3 — ítem inexistente:} el sistema informa que el ítem no existe en el catálogo y no registra la modificación.
\item \textbf{Paso 3 — ya existe un override para ese ítem:} el sistema reemplaza el anterior por el nuevo.
\item \textbf{Paso 5 — banear una categoría:} el DM elige una categoría (por ejemplo, magia); el sistema la marca como no disponible para toda la campaña.
\item \textbf{Quitar un override:} el DM elige un override y lo quita; el sistema devuelve la regla base.
\item \textbf{Catálogo externo no disponible:} el sistema informa el problema y muestra solo los overrides de la campaña; no se pueden validar ítems nuevos hasta que el catálogo responda.
\item \textbf{Efecto sobre personajes:} un ítem baneado se rechaza al crear o editar personajes y al repartir botín.
\item \textbf{Usuario sin permiso:} si el usuario no es DM de la campaña, el sistema rechaza la operación.
\end{itemize} \\ \hline
\end{longtable}}
\subparagraph{Prototipo de la interfaz de usuario}
Pantallas: RuleSet de la campaña (lista de reglas con marca de override, filtro por tipo); buscador del catálogo; formulario de modificación de estadísticas; acciones Banear y Quitar override.

\figdiagh{\subparagraph{Diagrama de casos de uso}}{regiardo_manuel/uc-cu-006-001}{Diagrama de casos de uso de CU-006-001.}{0.9}
\figdiagh{\subparagraph{Diagrama de robustez}}{regiardo_manuel/robustness-cu-006-001}{Diagrama de robustez de CU-006-001.}{0.9}
\figdiagh{\subparagraph{Diagrama de secuencia}}{regiardo_manuel/sequence-cu-006-001}{Diagrama de secuencia de CU-006-001.}{0.9}

% Generado por scripts/cap10_gen.py desde documentation/legacy. No editar a mano.
\paragraph{CU-002-001 -- Administrar Campañas}
{\small\setlength{\tabcolsep}{5pt}\renewcommand{\arraystretch}{1.2}
\begin{longtable}{|>{\raggedright\arraybackslash\bfseries}p{3.2cm}|>{\raggedright\arraybackslash}p{11.9cm}|}
\hline
Objetivo & Permitir al DM crear y administrar sus campañas: armar el grupo de jugadores, mantener los datos de la campaña y controlar quién puede entrar. \\ \hline
Identificación & CU-002-001 \\ \hline
Nombre & Administrar Campañas \\ \hline
Descripción & El usuario autenticado crea una campaña, elige su sistema de reglas base y pasa a ser su DM. El sistema genera un código de invitación único que el DM comparte con su grupo; los jugadores se unen con ese código. El DM consulta el detalle completo de la campaña, edita su nombre y descripción, regenera el código de invitación cuando quiere cerrar el acceso al anterior, remueve jugadores y, si lo desea, elimina la campaña. Los jugadores ven una vista reducida de la campaña, sin datos exclusivos del DM. \\ \hline
Precondición & \begin{itemize}[itemsep=2pt,topsep=2pt]
\item El usuario está autenticado.
\item Para editar, eliminar, gestionar el código o remover jugadores, el usuario es DM de la campaña.
\end{itemize} \\ \hline
Postcondición & \begin{itemize}[itemsep=2pt,topsep=2pt]
\item La campaña existe con sus datos actualizados, su código de invitación vigente y su lista de jugadores, o bien fue eliminada junto con sus datos dependientes.
\end{itemize} \\ \hline
Actores primarios & DM \\ \hline
Actores secundarios & Jugador (se une, consulta y abandona la campaña). \\ \hline
Eventos disparadores & El usuario quiere iniciar una campaña nueva o administrar una existente. \\ \hline
Extensiones y condiciones & \begin{itemize}[itemsep=2pt,topsep=2pt]
\item \textbf{include}: CU-002-002 Crear Campaña. Al iniciar una campaña (incluye CU-006-002 Asignar Sistema de Reglas Base)
\item \textbf{include}: CU-002-003 Consultar Campañas y Detalle. Siempre que se revisa una campaña
\item \textbf{extend}: CU-002-004 Editar Campaña. El DM cambia nombre o descripción
\item \textbf{extend}: CU-002-005 Eliminar Campaña. El DM decide eliminarla
\item \textbf{extend}: CU-002-006 Gestionar Código de Invitación. El DM consulta o regenera el código
\item \textbf{extend}: CU-002-007 Remover Jugador de la Campaña. El DM saca a un jugador
\item \textbf{extend}: CU-002-008 Unirse a una Campaña. Lo realiza el Jugador con un código válido
\item \textbf{extend}: CU-002-009 Abandonar una Campaña. Lo realiza el Jugador
\end{itemize}
Reglas: cada campaña tiene un único DM, su creador, que no puede removerse a sí mismo; el código de invitación es único.

Requisitos relacionados: RF-006 a RF-016, RD-001 a RD-003, RNF-001, RNF-002, RNF-004 (ver HU-01, HU-02, HU-04). \\ \hline
Escenario principal & \begin{enumerate}[itemsep=2pt,topsep=2pt]
\item \textbf{Usuario:} solicita crear una campaña e ingresa nombre, descripción y sistema de reglas.
\item \textbf{Sistema:} crea la campaña, lo designa DM y muestra el código de invitación.
\item \textbf{DM:} comparte el código con sus jugadores.
\item \textbf{Jugador:} ingresa el código de invitación.
\item \textbf{Sistema:} agrega al jugador a la campaña y la muestra en su lista.
\item \textbf{DM:} abre la campaña.
\item \textbf{Sistema:} muestra el detalle completo: datos, jugadores y código de invitación.
\item \textbf{DM:} modifica el nombre o la descripción.
\item \textbf{Sistema:} guarda los cambios y los muestra.
\end{enumerate} \\ \hline
Escenarios alternativos & \begin{itemize}[itemsep=2pt,topsep=2pt]
\item \textbf{Paso 1 — datos incompletos:} el sistema indica los campos faltantes y no crea la campaña.
\item \textbf{Paso 4 — código inválido:} el sistema informa que el código no existe y no agrega al jugador.
\item \textbf{Paso 4 — el jugador ya pertenece a la campaña:} el sistema informa que ya forma parte.
\item \textbf{Paso 7 — consulta de un Jugador:} el sistema muestra una vista reducida, sin datos exclusivos del DM.
\item \textbf{Regenerar código:} el DM solicita un código nuevo; el sistema invalida el anterior, muestra el nuevo y conserva a los jugadores ya unidos.
\item \textbf{Remover jugador:} el DM elige a un jugador; el sistema lo saca de la campaña. Si el DM intenta removerse a sí mismo, el sistema lo rechaza.
\item \textbf{Abandonar:} un Jugador abandona la campaña y el sistema lo saca de la lista.
\item \textbf{Eliminar campaña:} el DM confirma la eliminación; el sistema elimina la campaña con sus sesiones, personajes, RuleSet y bitácoras.
\item \textbf{Usuario sin permiso:} si el usuario no es DM de la campaña, el sistema rechaza editar, eliminar, regenerar el código o remover jugadores.
\end{itemize} \\ \hline
\end{longtable}}
\subparagraph{Prototipo de la interfaz de usuario}
Pantallas: Mis campañas (lista con rol); Crear campaña (nombre, descripción, sistema de reglas); Detalle de campaña (datos, jugadores, código de invitación con botón Regenerar, acciones Editar y Eliminar); Unirse con código.

\figdiagh{\subparagraph{Diagrama de casos de uso}}{regiardo_manuel/uc-cu-002-001}{Diagrama de casos de uso de CU-002-001.}{0.9}
\figdiagh{\subparagraph{Diagrama de robustez}}{regiardo_manuel/robustness-cu-002-001}{Diagrama de robustez de CU-002-001.}{0.9}
\figdiagh{\subparagraph{Diagrama de secuencia}}{regiardo_manuel/sequence-cu-002-001}{Diagrama de secuencia de CU-002-001.}{0.9}

% Generado por scripts/cap10_gen.py desde documentation/legacy. No editar a mano.
\paragraph{CU-005-001 -- Gestionar Eventos de Sesión}
{\small\setlength{\tabcolsep}{5pt}\renewcommand{\arraystretch}{1.2}
\begin{longtable}{|>{\raggedright\arraybackslash\bfseries}p{3.2cm}|>{\raggedright\arraybackslash}p{11.9cm}|}
\hline
Objetivo & Permitir al DM registrar durante la partida los hechos mecánicos que afectan a los personajes (daño, curación, experiencia, botín), de modo que el grupo los vea al instante y quede un historial consultable de lo ocurrido. \\ \hline
Identificación & CU-005-001 \\ \hline
Nombre & Gestionar Eventos de Sesión \\ \hline
Descripción & Con la sesión en curso, el DM registra un evento sobre un personaje, sobre varios o sobre todo el grupo: daño, curación, experiencia, botín o un suceso personalizado. El sistema aplica el efecto sobre los personajes, guarda el evento y lo difunde en tiempo real a los jugadores conectados, que ven su ficha actualizada sin recargar la página. Cada evento deja además una entrada automática en la bitácora de la sesión, que el DM puede completar con entradas de combate o generales. Todos los miembros de la campaña pueden consultar el historial de eventos y la bitácora. Un cliente que pierde la conexión recupera el estado al reconectarse. \\ \hline
Precondición & \begin{itemize}[itemsep=2pt,topsep=2pt]
\item El usuario está autenticado y es DM de la campaña.
\item La sesión está en curso.
\end{itemize} \\ \hline
Postcondición & \begin{itemize}[itemsep=2pt,topsep=2pt]
\item El evento queda registrado y sus efectos aplicados a los personajes destinatarios.
\item Existe una entrada de acción en la bitácora de la sesión.
\item Los jugadores conectados recibieron el evento.
\end{itemize} \\ \hline
Actores primarios & DM \\ \hline
Actores secundarios & Jugador (recibe los eventos en tiempo real y consulta el historial y la bitácora). \\ \hline
Eventos disparadores & Ocurre en la partida un hecho que afecta a uno o más personajes. \\ \hline
Extensiones y condiciones & \begin{itemize}[itemsep=2pt,topsep=2pt]
\item \textbf{include}: CU-005-002 Difundir Evento en Tiempo Real. Siempre, al registrar el evento
\item \textbf{include}: CU-005-003 Registrar Entrada de Acción en Bitácora. Siempre, al registrar el evento
\item \textbf{extend}: CU-005-004 Registrar Evento de Daño. El evento reduce los puntos de vida
\item \textbf{extend}: CU-005-005 Registrar Evento de Curación. El evento aumenta los puntos de vida
\item \textbf{extend}: CU-005-006 Registrar Evento de XP. El evento suma experiencia
\item \textbf{extend}: CU-005-007 Registrar Evento de Loot. El evento entrega ítems al personaje
\item \textbf{extend}: CU-005-008 Registrar Evento Personalizado. El evento solo deja constancia en el historial
\item \textbf{extend}: CU-005-009 Consultar Historial de Eventos. Lo realiza el DM o un Jugador
\item \textbf{extend}: CU-005-010 Reconstruir Estado tras Reconexión. Un cliente se reconecta
\item \textbf{extend}: CU-005-011 Gestionar Bitácora de Sesión. El DM registra entradas de combate o generales (incluye registrar entrada de combate, registrar entrada general y consultar bitácora)
\end{itemize}
Reglas: los puntos de vida del personaje se mantienen entre 0 y su máximo; un evento de loot con algún ítem baneado por el RuleSet de la campaña se rechaza; todos los miembros de la campaña ven todos los eventos de la sesión.

Requisitos relacionados: RF-035 a RF-041, RD-011, RD-012, RD-017 a RD-020, RNF-007, RS-004, RNF-001, RNF-002 (ver HU-10, HU-11). \\ \hline
Escenario principal & \begin{enumerate}[itemsep=2pt,topsep=2pt]
\item \textbf{DM:} elige el tipo de evento "daño", el personaje afectado e ingresa el valor 12.
\item \textbf{Sistema:} valida que la sesión esté en curso, aplica el daño al personaje y registra el evento.
\item \textbf{Sistema:} confirma al DM que el evento fue registrado.
\item \textbf{Jugador:} ve su ficha de personaje actualizada al instante, sin recargar.
\item \textbf{DM:} abre el historial de la sesión.
\item \textbf{Sistema:} muestra los eventos y las entradas de la bitácora en orden cronológico, incluida la entrada del daño recién registrado.
\end{enumerate} \\ \hline
Escenarios alternativos & \begin{itemize}[itemsep=2pt,topsep=2pt]
\item \textbf{Paso 1 — varios personajes o todo el grupo:} el DM elige a más de un personaje o a todo el grupo; el sistema aplica el evento a cada uno.
\item \textbf{Paso 1 — otros tipos de evento:} curación, experiencia, botín o personalizado; el sistema aplica el efecto propio de cada tipo. El evento personalizado no modifica estadísticas y solo deja constancia.
\item \textbf{Paso 2 — la sesión no está en curso:} el sistema rechaza el evento e informa que la sesión debe estar en curso.
\item \textbf{Paso 2 — el daño supera los puntos de vida:} el sistema deja los puntos de vida del personaje en 0.
\item \textbf{Paso 2 — la curación supera el máximo:} el sistema deja los puntos de vida en el máximo del personaje.
\item \textbf{Paso 2 — loot con un ítem baneado:} el sistema rechaza el evento e informa qué ítem está baneado en la campaña.
\item \textbf{Paso 4 — el jugador no está conectado:} no recibe el evento en vivo; al reconectarse, el cliente consulta la bitácora, el historial y su ficha y recupera el estado.
\item \textbf{Después del paso 3 — bitácora manual:} el DM registra una entrada de combate o general, opcionalmente asociada a personajes, y el sistema la agrega a la bitácora.
\item \textbf{Usuario sin permiso:} si el usuario no es DM de la campaña, el sistema rechaza el registro de eventos y de entradas de bitácora; la consulta del historial y de la bitácora está permitida a todos los miembros.
\end{itemize} \\ \hline
\end{longtable}}
\subparagraph{Prototipo de la interfaz de usuario}
Pantallas: Registrar evento (tipo, valor, personajes destinatarios o todo el grupo); Historial de la sesión (eventos y bitácora en orden cronológico, con filtro por tipo); Ficha del personaje con actualización en vivo (vista del jugador).

\figdiagh{\subparagraph{Diagrama de casos de uso}}{regiardo_manuel/uc-cu-005-001}{Diagrama de casos de uso de CU-005-001.}{0.9}
\figdiagh{\subparagraph{Diagrama de robustez}}{regiardo_manuel/robustness-cu-005-001}{Diagrama de robustez de CU-005-001.}{0.9}
\figdiagh{\subparagraph{Diagrama de secuencia}}{regiardo_manuel/sequence-cu-005-001}{Diagrama de secuencia de CU-005-001.}{0.9}

\fi

\subsubsection{Arquitectura}
% Generado por scripts/cap10_gen.py desde documentation/legacy. No editar a mano.
\paragraph{Vista General — 3 Capas Cliente-Servidor}
\begin{verbatim}
Cliente (Frontend)  ->  Servidor (Backend)  ->  Base de Datos (SQL Server)
\end{verbatim}
Además del backend principal, existe un \textbf{servicio externo independiente} para la generación de diálogo de NPC (ver sección 4), contactado por el backend pero desplegado y evolucionado por separado.

\paragraph{Frontend — Arquitectura MVC}
{\small\setlength{\tabcolsep}{4pt}\renewcommand{\arraystretch}{1.15}
\begin{longtable}{|>{\raggedright\arraybackslash}p{5.0cm}|>{\raggedright\arraybackslash}p{5.0cm}|>{\raggedright\arraybackslash}p{5.0cm}|}
\hline
\textbf{Capa} & \textbf{Responsabilidad} & \textbf{Implementación} \\ \hline\endhead
\textbf{View} & Presentación pura; renderiza el estado recibido. & \texttt{pages/\allowbreak{}}, \texttt{components/\allowbreak{}} (React) \\ \hline
\textbf{Controller} & Procesa eventos de entrada del usuario (submits, clicks, navegación) y los traduce en acciones sobre el Model. & Hooks y manejadores de eventos dentro de páginas/componentes \\ \hline
\textbf{Model} & Estado de la aplicación y acceso a datos remotos. & \texttt{redux/\allowbreak{}} (store, slices) + \texttt{api/\allowbreak{}} (cliente HTTP hacia el backend) \\ \hline
\end{longtable}}
El flujo es unidireccional: \textbf{View $\rightarrow$ Controller $\rightarrow$ Model $\rightarrow$ View} (el Model notifica cambios de estado que la View vuelve a renderizar), consistente con el diagrama general de arquitectura.

\paragraph{Backend — Sistema Estratificado Modularizado}
El backend sigue una \textbf{arquitectura estratificada (Stratified Systems)}, organizada en \textbf{módulos por dominio} (\texttt{auth}, \texttt{campaign}, \texttt{character}, \texttt{session}, \texttt{enemy}, \texttt{ruleset}, \texttt{npc-dialogue}, \texttt{event} y \texttt{session-log}). Cada módulo replica las mismas 4 capas:

\begin{verbatim}
Routes -> Controller -> Service -> Repository
\end{verbatim}
{\small\setlength{\tabcolsep}{4pt}\renewcommand{\arraystretch}{1.15}
\begin{longtable}{|>{\raggedright\arraybackslash}p{7.7cm}|>{\raggedright\arraybackslash}p{7.7cm}|}
\hline
\textbf{Capa} & \textbf{Responsabilidad} \\ \hline\endhead
\textbf{Routes} & Define los endpoints HTTP del módulo y aplica middleware (autenticación, autorización por rol de campaña, validación de esquema). \\ \hline
\textbf{Controller} & Adapta la request/response HTTP; delega la lógica al Service. Sin lógica de negocio. \\ \hline
\textbf{Service} & Lógica de negocio y de dominio (p. ej. generación de \texttt{inviteCode}, aplicación de \texttt{RuleOverride} vía Strategy + Decorator, orquestación de llamadas a servicios externos). \\ \hline
\textbf{Repository} & Encapsula el acceso a la persistencia (SQL Server). Desacopla la lógica de negocio del motor y la tecnología de persistencia concretos, y facilita testear los \texttt{Service} con un Repository simulado (mock/stub). \\ \hline
\end{longtable}}
Middleware transversal (\texttt{auth}, \texttt{userContext}, \texttt{campaignRole}, \texttt{validation}) se aplica a nivel de \texttt{Routes} antes de llegar al \texttt{Controller}.

\subparagraph{Entrega en tiempo real — WebSocket Gateway (módulo `session`)}
El módulo \texttt{session} incorpora, junto a las 4 capas estándar, un \textbf{WebSocket Gateway} (canal \texttt{session:\{id\}}). No es una capa adicional del stack Routes$\rightarrow$Controller$\rightarrow$Service$\rightarrow$Repository, sino un mecanismo de \textbf{entrega} paralelo al ciclo HTTP request/response:

\begin{itemize}[itemsep=2pt,topsep=2pt]
\item El \texttt{Service} de sesión, tras persistir un \texttt{SessionEvent} (y su reflejo automático en \texttt{SessionLog}, ver CU-005-001), emite el evento por WebSocket a los clientes conectados al canal de esa sesión.
\item Los jugadores afectados ven su ficha de personaje actualizarse sin recargar la página.
\item \textbf{SQL Server sigue siendo la única fuente de verdad.} El WebSocket es puramente de entrega en tiempo real: si un cliente se desconecta y reconecta, reconstruye el estado consultando la bitácora (CU-005-011) y la ficha del personaje vía REST, no depende de haber recibido el mensaje de socket.
\end{itemize}
\paragraph{Servicio de Generación de Diálogo de NPC — Pipes \& Filters + Blackboard}
Decisión de \textbf{arquitectura}: el servicio se implementa como un \textbf{servicio externo independiente}, no como un módulo más del backend estratificado.

\subparagraph{Problema}
\begin{itemize}[itemsep=2pt,topsep=2pt]
\item El DM necesita, en medio de una sesión, generar diálogo genérico para un NPC no planificado, sin bloquear el resto del sistema ante la latencia o eventuales fallos del proveedor de IA.
\item Armar un buen prompt requiere combinar información heterogénea y de distintas fuentes (tono/ambientación de la campaña, ficha del NPC, contexto de la escena, reglas activas) antes de poder invocar al modelo de lenguaje.
\item El servicio tiene un ciclo de evolución y de escalado propio (agregar/quitar filtros, cambiar de proveedor de IA) y no debe acoplar la disponibilidad del backend principal a la latencia de un proveedor externo.
\end{itemize}
\subparagraph{Solución}
\begin{itemize}[itemsep=2pt,topsep=2pt]
\item El servicio se estructura como \textbf{Pipes \& Filters}: una secuencia de filtros independientes conectados por "pipes", donde la salida de un filtro es la entrada del siguiente (recibir solicitud y contexto, construir el prompt, invocar el modelo de IA, formatear y validar la respuesta).
\item El filtro de construcción del prompt se resuelve internamente con una arquitectura \textbf{Blackboard}: múltiples "fuentes de conocimiento" (contexto de campaña, perfil del NPC, tono narrativo, reglas activas de la campaña) escriben sus aportes sobre una pizarra compartida; un controlador decide cuándo la información acumulada es suficiente para ensamblar el prompt final.
\item El filtro que invoca al modelo oculta al proveedor de IA concreto detrás de una \textbf{Facade} (patrón de diseño; ver \texttt{design-patterns.\allowbreak{}md}, sección 2).
\item El backend lo consume desde el módulo \texttt{npc-dialogue} (capa \texttt{Service}) como un cliente HTTP, análogo a la integración con la API DnD5e: le envía el contexto necesario y recibe el texto generado, sin conocer los detalles internos del pipeline.
\item La persistencia del resultado es \textbf{opcional y manual}: el DM decide si el diálogo amerita guardarse como una entrada de \texttt{SessionLog} (\texttt{type = general}); no ocurre automáticamente.
\end{itemize}
\subparagraph{Por qué esta arquitectura}
\begin{itemize}[itemsep=2pt,topsep=2pt]
\item \textbf{Pipes \& Filters} permite evolucionar el pipeline (agregar un filtro de moderación, cachear resultados, cambiar de proveedor de IA) sin afectar al resto de los filtros, y facilita el despliegue y escalado independientes respecto del backend principal.
\item \textbf{Blackboard} es apropiado porque construir el prompt no es un pipeline lineal simple: distintas fuentes de contexto (campaña, NPC, reglas, escena) contribuyen de forma no necesariamente secuencial y con prioridades variables; el Blackboard modela ese ensamblado colaborativo antes de "cerrar" el prompt.
\end{itemize}
\paragraph{Decisiones Tomadas}
\begin{itemize}[itemsep=2pt,topsep=2pt]
\item \textbf{Autenticación}: propia, basada en \textbf{JWT} (registro y login con email + password persistidos en SQL Server). No se depende de un proveedor de identidad externo.
\item \textbf{Persistencia}: \textbf{SQL Server} con \textbf{Prisma} como ORM y sistema de migraciones; el acceso a datos queda encapsulado en la capa \texttt{Repository} de cada módulo.
\item \textbf{Servicio de diálogo de NPC}: servicio externo contactado por \textbf{HTTP/JSON} desde el módulo \texttt{npc-dialogue}.
\item \textbf{Tiempo real}: \textbf{WebSocket} (canal \texttt{session:\{id\}}) como mecanismo de entrega.
\end{itemize}

\figdiag{regiardo_manuel/component-architecture}{Arquitectura general del sistema.}{0.9}
\figdiag{regiardo_manuel/component-npc-dialogue-service}{Servicio de generación de diálogo de NPC (Pipes \& Filters con Blackboard).}{0.9}

\subsubsection{Patrones de diseño}
% Generado por scripts/cap10_gen.py desde documentation/legacy. No editar a mano.
\paragraph{Strategy — Asignación de RuleSets}

\subparagraph{Problema}
\begin{itemize}[itemsep=2pt,topsep=2pt]
\item El sistema debe soportar múltiples sistemas de reglas de juego de mesa (D\&D 5e hoy, potencialmente Pathfinder u otros a futuro), sin que agregar un sistema nuevo requiera modificar el código existente.
\end{itemize}
\subparagraph{Solución}
\begin{itemize}[itemsep=2pt,topsep=2pt]
\item \textbf{Strategy} (\texttt{IRuleSetProvider}): cada sistema de reglas (D\&D 5e, Pathfinder, ...) implementa una interfaz común para resolver items, hechizos, armas y monstruos. Agregar un sistema nuevo implica únicamente escribir una nueva clase concreta (\texttt{DnD5eRuleSetProvider}, \texttt{PathfinderRuleSetProvider}), respetando el principio Open/Closed.
\item Una \textbf{Factory} (\texttt{RuleSetProviderFactory}) resuelve qué Strategy concreta instanciar según el \texttt{RuleSystem} asignado a la campaña (\texttt{CampaignRuleSet.\allowbreak{}rule\_\allowbreak{}system\_\allowbreak{}id}).
\end{itemize}
\subparagraph{Por qué Strategy y no Template Method}
\begin{itemize}[itemsep=2pt,topsep=2pt]
\item \textbf{Template Method} no aplica: no hay un algoritmo con pasos fijos y "huecos" a redefinir por subclases; lo que varía es el \textit{origen de los datos} (sistema de reglas) y una \textit{capa opcional de reemplazo} (personalización), que es exactamente el par Strategy/Decorator.
\end{itemize}
\subparagraph{Comparación canónica}
{\small\setlength{\tabcolsep}{4pt}\renewcommand{\arraystretch}{1.15}
\begin{longtable}{|>{\raggedright\arraybackslash}p{3.4cm}|>{\raggedright\arraybackslash}p{11.9cm}|}
\hline
\textbf{Título} & \textbf{Contenido} \\ \hline\endhead
\textbf{Nombre de Implementación} & Asignación de RuleSets \\ \hline
\textbf{Patrón} & Strategy \\ \hline
\textbf{Cambios previstos} & Se sabe de antemano que se agregarán nuevos sistemas de reglas (Pathfinder, etc.) sin tocar código existente (Open/Closed). \\ \hline
\textbf{Funcionalidad} & Cada sistema de reglas resuelve items/hechizos/armas/monstruos con lógica propia (p. ej. distinta API o catálogo fuente). \\ \hline
\textbf{Restricciones de Diseño} & El llamador (\texttt{CampaignRuleSetDecorator}) debe poder tratar cualquier sistema de reglas de forma intercambiable y uniforme en tiempo de ejecución, seleccionado dinámicamente por campaña (\texttt{CampaignRuleSet.\allowbreak{}rule\_\allowbreak{}system\_\allowbreak{}id}). \\ \hline
\textbf{Dónde} & \textit{\texttt{Strategy}} \textbf{es} \textit{\texttt{IRuleSetProvider}} \newline{} \textit{\texttt{ConcreteStrategyA}} \textbf{es} \textit{\texttt{DnD5eRuleSetProvider}} \newline{} \textit{\texttt{ConcreteStrategyB (futuro)}} \textbf{es} \textit{\texttt{PathfinderRuleSetProvider}} \newline{} \textit{\texttt{Context}} \textbf{es} \textit{\texttt{CampaignRuleSetDecorator}} (mantiene la referencia \texttt{wrapped: IRuleSetProvider}) \newline{} \textit{\texttt{Selección de Strategy}} \textbf{es} \textit{\texttt{RuleSetProviderFactory.\allowbreak{}createFor(ruleSystem)}} \\ \hline
\textbf{Comentarios} & Clave: el Strategy resuelve el catálogo \textbf{base} (sin overrides); la personalización por campaña es responsabilidad exclusiva del Decorator (tabla siguiente), nunca del Strategy. \\ \hline
\end{longtable}}
\figdiagh{}{regiardo_manuel/class-ruleset-patterns}{Patrones Strategy, Decorator y Factory para los RuleSets.}{0.9}
\paragraph{Decorator — Personalización de RuleSet por Campaña}

\subparagraph{Problema}
\begin{itemize}[itemsep=2pt,topsep=2pt]
\item Dentro de una campaña puntual, el DM debe poder personalizar ese sistema base: banear armas/hechizos o modificar estadísticas de un item, sin alterar el catálogo compartido por el resto de las campañas.
\end{itemize}
\subparagraph{Solución}
\begin{itemize}[itemsep=2pt,topsep=2pt]
\item \textbf{Decorator} (\texttt{CampaignRuleSetDecorator}): envuelve al Strategy base seleccionado por la campaña. Ante cada consulta, revisa primero si existe un \texttt{RuleOverride} aplicable (ban / modify\_stat); si lo hay, lo aplica (o corta la respuesta indicando que el ítem está baneado); si no, delega transparentemente en el proveedor base envuelto.
\end{itemize}
\subparagraph{Por qué Decorator y no Chain of Responsibility}
\begin{itemize}[itemsep=2pt,topsep=2pt]
\item \textbf{Chain of Responsibility} para los overrides fue considerado, pero un ítem solo tiene, como máximo, un override relevante por campaña (ban \textit{o} modify\_stat) — no hay una cadena de handlers que se pasen la solicitud entre sí, sino una única consulta indexada por \texttt{target\_\allowbreak{}key}. Decorator expresa mejor la idea de "capa que envuelve y opcionalmente reemplaza el comportamiento base".
\end{itemize}
\subparagraph{Comparación canónica}
{\small\setlength{\tabcolsep}{4pt}\renewcommand{\arraystretch}{1.15}
\begin{longtable}{|>{\raggedright\arraybackslash}p{3.4cm}|>{\raggedright\arraybackslash}p{11.9cm}|}
\hline
\textbf{Título} & \textbf{Contenido} \\ \hline\endhead
\textbf{Nombre de Implementación} & Personalización de RuleSet por Campaña \\ \hline
\textbf{Patrón} & Decorator \\ \hline
\textbf{Cambios previstos} & El conjunto de tipos de override (\texttt{ban}, \texttt{modify\_\allowbreak{}stat}, y otros a futuro como \texttt{modify\_\allowbreak{}cost}) puede crecer sin alterar los proveedores base ni el resto del sistema. \\ \hline
\textbf{Funcionalidad} & Un DM puede banear o modificar estadísticas de items/hechizos solo dentro de su campaña, sin afectar el catálogo compartido por otras campañas que usan el mismo \texttt{RuleSystem}. \\ \hline
\textbf{Restricciones de Diseño} & La personalización debe ser \textbf{opcional y componible} sobre cualquier Strategy base, sin que el Strategy sepa que está siendo envuelto (principio de responsabilidad única). \\ \hline
\textbf{Dónde} & \textit{\texttt{Component}} \textbf{es} \textit{\texttt{IRuleSetProvider}} (misma interfaz común que Strategy) \newline{} \textit{\texttt{ConcreteComponent}} \textbf{es} \textit{\texttt{DnD5eRuleSetProvider}} / \textit{\texttt{PathfinderRuleSetProvider}} \newline{} \textit{\texttt{Decorator}} \textbf{es} \textit{\texttt{CampaignRuleSetDecorator}} (atributo \texttt{wrapped: IRuleSetProvider}, agrega \texttt{overrides: RuleOverride[]}) \\ \hline
\textbf{Comentarios} & Clave: se descartó \textbf{Chain of Responsibility} porque no hay una cadena de handlers — cada consulta tiene, como máximo, \textbf{un} override aplicable (búsqueda indexada por \texttt{target\_\allowbreak{}key}), lo cual encaja mejor con "envolver y opcionalmente reemplazar" (Decorator) que con "pasar la solicitud por una cadena". \\ \hline
\end{longtable}}
\paragraph{Facade — Selección de Proveedor de IA}

\subparagraph{Problema}
\begin{itemize}[itemsep=2pt,topsep=2pt]
\item El filtro "Invocar Modelo de IA" del pipeline del servicio de diálogo (ver la sección de arquitectura) necesita comunicarse con un proveedor de IA concreto (OpenAI, un modelo auto-hospedado, u otro a futuro), cada uno con su propio esquema de autenticación, formato de request/response, límites de tokens y manejo de errores.
\item No queremos que el resto del pipeline (ni el backend principal) conozca esos detalles, ni que cambiar de proveedor implique tocar código disperso por varios puntos del sistema.
\end{itemize}
\subparagraph{Solución}
\begin{itemize}[itemsep=2pt,topsep=2pt]
\item \textbf{Facade} (\texttt{AiProviderFacade}): expone una única operación simple (\texttt{generateCompletion(prompt, options)}) al filtro que la consume, ocultando toda la complejidad de los subsistemas concretos (\texttt{OpenAiClient}, \texttt{LocalLlmClient}, ...).
\item El Facade resuelve internamente, a partir de una configuración externa (\texttt{AiProviderConfig.\allowbreak{}activeProvider}), a cuál cliente concreto delegar la llamada. Agregar un nuevo proveedor implica escribir un nuevo cliente que implemente \texttt{ILlmClient}, sin modificar el Facade ni el filtro que lo invoca.
\item A diferencia del par Strategy + Decorator (secciones 1 y 2), aquí \textbf{no} hay necesidad de una interfaz de estrategia expuesta al llamador ni de una capa de "override": el objetivo es únicamente \textbf{simplificar y desacoplar} el acceso a un subsistema externo complejo y cambiante, que es exactamente el propósito de Facade.
\end{itemize}
\subparagraph{Por qué Facade y no Strategy}
\begin{itemize}[itemsep=2pt,topsep=2pt]
\item Se consideró reutilizar Strategy (como en RuleSet) para el proveedor de IA, ya que ambos casos "seleccionan una implementación concreta según configuración". La diferencia clave es que en RuleSet el llamador (\texttt{CampaignRuleSetDecorator}) necesita tratar a la Strategy como \textbf{intercambiable en tiempo de ejecución por campaña} y además envolverla con una capa de personalización (Decorator) — hay variabilidad de negocio real que el llamador debe modelar. En cambio, el llamador del proveedor de IA (el filtro "Invocar Modelo de IA") no necesita saber que existen múltiples proveedores ni razonar sobre cuál usar: solo quiere "un texto a partir de un prompt". Facade es más simple y ajustado quirúrgicamente a ese caso: \textbf{una interfaz mínima que oculta subsistemas heterogéneos}, sin exponer la variabilidad como parte del contrato.
\end{itemize}
\subparagraph{Comparación canónica}
{\small\setlength{\tabcolsep}{4pt}\renewcommand{\arraystretch}{1.15}
\begin{longtable}{|>{\raggedright\arraybackslash}p{3.4cm}|>{\raggedright\arraybackslash}p{11.9cm}|}
\hline
\textbf{Título} & \textbf{Contenido} \\ \hline\endhead
\textbf{Nombre de Implementación} & Selección de Proveedor de IA \\ \hline
\textbf{Patrón} & Facade \\ \hline
\textbf{Cambios previstos} & El proveedor de IA (OpenAI, modelo local auto-hospedado, u otro) puede cambiarse por configuración o evolucionar sin que el pipeline que lo invoca se entere. \\ \hline
\textbf{Funcionalidad} & Cada proveedor tiene su propio esquema de autenticación, formato de request/response y manejo de límites/errores — el llamador solo necesita "generar texto a partir de un prompt". \\ \hline
\textbf{Restricciones de Diseño} & El filtro "Invocar Modelo de IA" del pipeline (Pipes \& Filters) no debe conocer subsistemas externos heterogéneos; a diferencia de RuleSet, aquí el llamador \textbf{no necesita razonar sobre cuál proveedor está activo}, solo consumir una interfaz mínima y estable. \\ \hline
\textbf{Dónde} & \textit{\texttt{Facade}} \textbf{es} \textit{\texttt{AiProviderFacade}} (método \texttt{generateCompletion(prompt, options)}) \newline{} \textit{\texttt{Subsistemas complejos}} \textbf{son} \textit{\texttt{OpenAiClient}}, \textit{\texttt{LocalLlmClient}} (futuro), ambos detrás de \texttt{ILlmClient} \newline{} \textit{\texttt{Configuración de selección}} \textbf{es} \textit{\texttt{AiProviderConfig.\allowbreak{}activeProvider}} \\ \hline
\textbf{Comentarios} & Clave: \textbf{no se usó Strategy} aquí porque el llamador no necesita tratar la variabilidad como parte de su contrato — solo desea ocultar la complejidad de un subsistema externo, que es el propósito específico de Facade. \\ \hline
\end{longtable}}
\figdiagh{}{regiardo_manuel/class-ai-provider-facade}{Patrón Facade para la selección del proveedor de IA.}{0.9}


\figdiagh{\subsubsection{Diagrama de clases}
Se presenta la vista general del backend y, luego, el diagrama de clases de cada caso de uso núcleo.}{regiardo_manuel/class-full}{Diagrama de clases general del backend.}{0.9}
\figdiag{regiardo_manuel/class-cu-004-001}{Diagrama de clases de CU-004-001.}{0.9}
\figdiag{regiardo_manuel/class-cu-007-001}{Diagrama de clases de CU-007-001.}{0.9}
\iftfi
\figdiag{regiardo_manuel/class-cu-006-001}{Diagrama de clases de CU-006-001.}{0.9}
\figdiag{regiardo_manuel/class-cu-002-001}{Diagrama de clases de CU-002-001.}{0.9}
\figdiag{regiardo_manuel/class-cu-005-001}{Diagrama de clases de CU-005-001.}{0.9}
\fi

\subsubsection{Especificación de casos de prueba}
\pendtxt
